\Changel
\section*{Introduction}

From Matsushima's formula and computations of $(\mathfrak G,K_\oo)$-cohomology,
we know that tempered automorphic representations contributions in the cohomology
of Shimura varieties with complex coefficients, is concentrated in middle degree.
If~one considers cohomology with coefficients in a very regular local system, then
only tempered representations can contribute so that all of the cohomology is concentrated
in middle degree.

For $\overline \Zm_l$-coefficients and Shimura varieties of Kottwitz-Harris-Taylor types,
we proved in \cite{boyer-stabilization}, whatever the weight of
the coefficients is, when the level is large enough at $l$, there are always nontrivial torsion
cohomology classes, so that the $\overline \Fm_l$\nobreakdash-co\-ho\-mology can not be concentrated in middle
degree. Thus if one wants an $\overline \Fm_l$-analog of the previous $\overline \Qm_l$-statement,
one must cut off some part of the cohomology.

In \cite{boyer-imj} for KHT Shimura varieties, and more generally
in \cite{scholze-cara} for any PEL proper Shimura variety, we obtain such a result
under some genericity hypothesis which can be stated as follows.
Let $(\sh_K)_{K \subset G(\Am^\oo)}$ be a tower, indexed by open compact subgroups
$K$ of $G(\Am^\oo)$, of compact Shimura varieties of Kottwitz type associated to some
similitude group $G$ and of relative dimension $d-1$ over its reflex field $F=EF^+$
where $F^+$ is totally real and $E/\Qm$ is an imaginary quadratic extension.
Let then $\mathfrak m$ be a system of Hecke eigenvalues appearing in
$H^{n_0}(\sh_K \times_F \overline F,\overline \Fm_l)$.
By the main result of \cite{scholze-torsion}, one can attach to such an $\mathfrak m$,
a mod $l$ Galois representation
\[
\overline{\rho_{\mathfrak m}}: \gal(\overline F/F) \to \GL_d(\overline \Fm_l).
\]
From \cite[Def.\,19]{scholze-cara},
we say that $\mathfrak m$ is generic (\resp decomposed generic) at some split $p$ in $E$,
if for all places $v$ of $F$ dividing $p$,
the set $\{\lambda_1,\dots,\lambda_d\}$ of eigenvalues of
$\overline{\rho_{\mathfrak m}}(\frob_v)$
satisfies $\lambda_i/\lambda_j \not \in \{q_v^{\pm 1}\}$ for all $i \neq j$
(\resp and are pairwise distinct), where $q_v$ is the cardinal
of the residue field at $v$. Then under the hypothesis\footnote{In their new preprint, Caraiani and Scholze explained that, from an observation of Koshikawa, one can replace decomposed generic
by simply generic, in their main statement.} that there exists such a $p$ with
$\mathfrak m$ generic at $p$, the integer $n_0$ above is necessarily equal to
the relative dimension of $\sh_K$. In particular the
$H^i(\sh_K \times_F \overline F,\overline \Zm_l)_{\mathfrak m}$
are all torsion-free.

In this work we consider the particular case of Kottwitz-Harris-Taylor Shimura varieties $\sh_K$
of \cite{h-t} associated to inner forms of $\GL_d$. Exploiting the fact, which is particular to these Shimura
varieties, that the non supersingular Newton strata are geometrically induced, we are then able to
prove the following result which happens to be useful at least for our approach of
Ihara's lemma, \cf \cite{boyer-ihara}.

\begin{thm*}
We suppose that $[F(\exp(2i\pi/l):F]>d$.
Let $\mathfrak m$ be a system of Hecke eigenvalues
such that $\overline{\rho_{\mathfrak m}}$
is irreducible, then the localized cohomology groups of $\sh_K$ with coefficients in any
$\overline \Zm_l$-local system $V_\xi$, are all free.
\end{thm*}

\skpt
\begin{remarks}
\begin{itemize}
\item
From \cite[\S4]{boyer-compositio}, we know that outside
the middle degree the irreducible Galois subquotients of
$H^i(\sh_K \times_F \overline F,\overline \Zm_l)_{\mathfrak m} \otimes_{\overline \Zm_l} \overline \Qm_l$ are of dimension strictly less than $d$, so~that
over $\overline \Qm_l$, the cohomology, localized at $\mathfrak m$,
is concentrated in middle degree. The previous theorem then tells us
it is the same for the $\overline \Fm_l$-cohomology.

\item
Note that Koshikawa, \cf \cite{koshikawa}, starting from \cite{boyer-imj} and using techniques from group theory, proved a similar result in low dimension.
\end{itemize}
\end{remarks}

\begin{lem*}
There exist infinitely many places $v$ of $F$ such that $q_v$, the order of the residue field of $F$ at $v$,
is of order strictly greater than $d$ in $(\Zm/l\Zm)^\times$.
\end{lem*}

\begin{proof}
The hypothesis $[F(\exp(2i\pi/l):F]>d$ is equivalent to the fact that the minimal polynomial
of $\exp(2i\pi/l)$ over $F$ is of degree $\delta> d$. By Cebotarev's
theorem, there exists then a set of places $v$ of strictly positive density such that the minimal polynomial of $\exp(2i\pi/l)$ over the residue field at $v$
is also of degree $\delta$. Recall that the roots
of this minimal polynomial are $\exp(2i\pi/l)^{q_v^k}$ with $k=0,\dots,
\delta-1$ with $\exp(2i\pi/l)^{q_v^{\delta}} =1$ in the residue field.
We then deduce that the order of $q_v$ modulo~$l$ divides $\delta$ but,
as $\exp(2i\pi/l)^{q_v^k} \neq 1$ in the residue field for $0 \leq k < \delta$,
then $q_v$ is of order $\delta > d$ modulo $l$.
\end{proof}

This property is used at three places in the proof.
\begin{itemize}
\item For a place $v$ as above, there is no irreducible cuspidal representation
$\pi_v$ of $\GL_g(F_v)$ with $g>1$ such that its reduction modulo $l$ has a supercuspidal
support made of characters, \cf the remark after Notation \ref{nota-mvarrho}.
This simplification is completely harmless and if one wants to take care of
these cuspidal representations, it~suffices to use \cite[Prop.\,2.4.1]{boyer-duke}.

\item With this hypothesis we also note that the pro-order of $\GL_d(\OC_v)$ is invertible
modulo $l$ so that, concerning torsion cohomology classes,
we can easily pass from infinite to maximal level at $v$, \cf for example Lemma \ref{lem-iIv}.

\item Finally in the last section, arguing by contradiction,
we are able to construct a sequence of intervals
$\{\lambda, \lambda q_v,\dots, \lambda q_v^{r}\}$ contained in the set of eigenvalues of
$\overline \rho_{\mathfrak m}(\frob_v)$ so that at the end we obtain the full set
$\{\lambda q_v^n \in \overline \Fm_l\sep n \in \Zm\}$ which is of order the order of $q_v$ modulo $l$, which
is trivially absurd if this order is strictly greater than the dimension $d$ of
$\overline \rho_{\mathfrak m}$.
\end{itemize}

The proof takes place in four main steps.

\begin{enumerate}
\item
First, in Section \ref{para-torsion1}, we analyze the torsion in the
cohomology of Harris-Taylor perverse sheaves at some place
$v$ and with infinite level at $v$, and we focus on $l$-torsion cohomology
classes with maximal non degeneracy level, \cf Lemma \ref{lem-rem-after1}.

\item
Recall, \cf \cite{boyer-torsion}, that one can define an exhaustive filtration
of stratification $\Fill^{\csbullet}(\Psi_v)$ of the nearby cycles perverse sheaf
whose graded parts are perverse sheaves
\[
\lexp p j^{=tg}_{!*} HT(\pi_v,\st_t(\pi_v)) \htarrowp P(t,\pi_v) \htarrowp
\lexp {p+} j^{=tg}_{!*} HT(\pi_v,\st_t(\pi_v)),
\]
\cf Section \ref{para-HT} for more details about $p$ and $p+$ perverse structures
and bimorphisms. The main point is that the constructions of \cite{boyer-torsion}
is of geometric nature and so works whatever the coefficients are.

Then on can
compute the cohomology of the Shimura variety through the
spectral sequence associated to the filtration
$\Fill^{\csbullet}(\Psi_v)$, whose
$E_1$ terms are given by the cohomology groups of the
Harris-Taylor perverse $\overline \Zm_l$-sheaves.
The main point, \cf Lemma~\ref{lem-generic-sq}, is that the $l$-torsion of the
cohomology of the Shimura variety with infinite level at $v$ does not have
an irreducible subquotient whose cuspidal support is made of characters.

\item
We then consider levels which are of Iwahori type at $v$ and infinite
at a place~$w$ verifying the same
hypothesis as $v$ above. We can then resume the previous result for~$w$, \ie the cohomology of the Shimura variety with infinite level at $w$ does not have
an irreducible subquotient whose cuspidal support is made of characters.
We~then focus on the various
lattices of $H^{d-1}_{\free}(\sh_{K} \times_F \overline F,V_{\xi,\overline \Qm_l})_{\widetilde{\mathfrak m}}$ where $\widetilde{\mathfrak m} \subset
\mathfrak m$ can be seen as a near equivalence class $\Pi_{\widetilde{\mathfrak m}}$
in the sense of of \cite{y-t}. The ones given by the
$\overline \Zm_l$\nobreakdash-cohomology,
are only slightly modified from the ones given by the cohomology of the
Harris-Taylor perverse sheaves, in the sense that the $l$-torsion of the
cokernel measuring the difference between two such lattices, as a representation
of $\GL_d(F_w)$, does not have any
irreducible generic subquotient with cuspidal
support made of characters, \cf Proposition \ref{prop-lattice-psi}.
Moreover if the torsion sub-module of
$H^{d-1}(\sh_{K} \times_F \overline F,V_{\xi,\overline \Zm_l})_{\mathfrak m}$ were not trivial, we prove, using the geometrical induced structure of the
Newton strata, that it would exist $\widetilde{\mathfrak m}$
such that the lattices mentioned above, are not isomorphic, \cf Proposition~\ref{prop-lattice-noniso}.

\item
Finally in Section \ref{para-final},
for any $\widetilde{\mathfrak m} \subset \mathfrak m$,
$H^{d-1}(\sh_K \times_F \overline F,V_{\xi,\overline \Zm_l})_{\mathfrak m}$
induces a quotient stable lattice $\Gamma_{\widetilde{\mathfrak m}}$ of $(\Pi_{\widetilde{\mathfrak m}}^{\oo})^K \otimes
\rho_{\widetilde{\mathfrak m}}$. As
$\overline \rho_{\mathfrak m}$ is supposed to be irreducible, then this
lattice is isomorphic to a tensor product of a stable lattice of
$(\Pi_{\widetilde{\mathfrak m}}^{\oo})^K$ by a stable lattice of
$\rho_{\widetilde{\mathfrak m}}$. Then the idea is to start from the filtration
of the free quotient of $H^{d-1}(\sh_K \times_F \overline F,
V_{\xi,\overline \Zm_l})_{\mathfrak m}$ given
by the filtration of the nearby cycles perverse sheaf, so~that, using repeatedly
diagrams as \eqref{eq-prop-extension}, we arrive at $\Gamma_{\widetilde{\mathfrak m}}$.
In the process we are able to construct an increasing sequence of intervals
contained in the set of eigenvalues of $\overline \rho_{\mathfrak m}(\frob v)$
so that at the end we obtain a full set $\{\lambda q_v^n \sep n \in \Zm\}$
which is of order the order of $q_v$ modulo $l$ which is, by hypothesis,
strictly greater than the dimension of $\overline \rho_{\mathfrak m}$, which is
absurd.
\end{enumerate}

We refer the reader to the introduction of Section \ref{para-proof} for more details.

\section{\texorpdfstring{Reminder from \cite{boyer-imj}}{R}}

\subsection{Representations of $\GL_d(K)$}
\label{para-gen}

Let $K/\Qm_p$ be a finite extension with $\OC_K$ its ring of integers,
$\varpi$ an uniformizer, and
$\kappa$ its residue field with order $q$.
We denote by $|\csbullet|$ its absolute value.
For a representation $\pi$ of $\GL_d(K)$ and $n \in \frac{1}{2} \Zm$, set
\[
\pi \{n\}:= \pi \otimes q^{-n \val \circ \det}.
\]
The Zelevinsky line associated to $\pi$ is by definition the set
$\{\pi \{n\}\sep n \in \Zm\}$.

\begin{notas} \label{nota-ind}
For $\pi_1$ and $\pi_2$ representations of respectively $\GL_{n_1}(K)$ and
$\GL_{n_2}(K)$, we will denote by
\[
\pi_1 \times \pi_2:=\ind_{P_{n_1,n_1+n_2}(K)}^{\GL_{n_1+n_2}(K)}
\pi_1 \{\sfrac{n_2}{2}\} \otimes \pi_2 \{-\sfrac{n_1}{2}\},
\]
the normalized parabolic induced representation where for any sequence
\[
\underline r=(0< r_1 < r_2 < \cdots < r_k=d),
\]
we write $P_{\underline r}$ for
the standard parabolic subgroup of $\GL_d$ with Levi
\[
\GL_{r_1} \times \GL_{r_2-r_1} \times \cdots \times \GL_{r_k-r_{k-1}}.
\]
\end{notas}

Recall that a representation
$\varrho$ of $\GL_d(K)$ is called \emph{cuspidal} (\resp \emph{supercuspidal})
if~it is not a subspace (\resp subquotient) of a proper parabolic induced representation.
When the field of coefficients is of characteristic zero then these two notions coincides,
but this is no more true for $\overline \Fm_l$.

\begin{defin}[{see \cite[\S9]{zelevinski2} and \cite[\S1.4]{boyer-compositio}}] \label{defi-rep}
Let $g$ be a divisor of $d=sg$ and $\pi$ an irreducible cuspidal
$\overline \Qm_l$-representation of $\GL_g(K)$.
The induced representation
\[
\pi\{\psfrac{1-s}{2}\} \times \pi \{\psfrac{3-s}{2}\} \times \cdots \times \pi \{\psfrac{s-1}{2}\}
\]
admits a unique irreducible quotient (\resp subspace) denoted $\st_s(\pi)$ (\resp $\speh_s(\pi)$); it is a generalized Steinberg (\resp Speh) representation.

Moreover the induced representation $\st_t(\pi \{\sfrac{-r}{2}\}) \times \speh_r(\pi \{\sfrac{t}{2}\})$
(\resp of $\st_{t-1}(\pi \{\psfrac{-r-1}{2}\}) \times \speh_{r+1}(\pi \{\psfrac{t-1}{2}\})$)
has a unique irreducible subspace (\resp quotient), denoted
$LT_\pi(t-1,r)$.
\end{defin}

\begin{remark}
These representations $LT_\pi(t-1,r)$ appear in the cohomology of the Lubin-Tate
spaces, \cf \cite{boyer-invent2}.
\end{remark}

\begin{prop}[{\cf \cite[III.5.10]{vigneras-livre}}] \label{prop-red-modl}
Let $\pi$ be an irreducible cuspidal representation of $\GL_g(K)$ with a stable
$\overline \Zm$-lattice\footnote{We say that $\pi$ is integral.}, then its reduction modulo $l$
is irreducible and cuspidal but not necessary supercuspidal.
\end{prop}

The supercuspidal support of the reduction modulo $l$ of a cuspidal representation,
is a segment associated to some irreducible $\overline \Fm_l$-supercuspidal representation
$\varrho$ of $\GL_{g_{-1}(\varrho)}(F_v)$ with $g=g_{-1}(\varrho) t$, where $t$ is
either equal to $1$ or of the following shape
$t=m(\varrho)l^u$ with $u \geq 0$ and where $m(\varrho)$ is defined as follows.

\begin{nota} \label{nota-mvarrho}
We denote by $m(\varrho)$ the order of the Zelevinsky line
$\{\varrho(\delta)\sep \delta \in \Zm\}$ of
$\varrho$ if it is not equal to $1$, otherwise $m(\varrho)=l$.
\end{nota}

\begin{remark}
When $\varrho$ is the trivial representation then $m(1_v)$ is either the order of $q$ modulo $l$
when it is $>1$, otherwise $m(1_v)=l$.
We say that such a $\pi_v$ is of $\varrho$-type~$u$ with $u \geq -1$.
\end{remark}


\begin{nota}
For $\varrho$ an irreducible $\overline \Fm_l$-supercuspidal representation, we denote by
$\cusp_\varrho$ (\resp $\cusp_\varrho(u)$ for some $u \geq -1$)
the set of equivalence classes of irreducible
$\overline \Qm_l$-cuspidal representations whose reduction modulo $l$ has
for supercuspidal
support a segment associated to $\varrho$ (\resp of $\varrho$-type $u$).
\end{nota}

Let $u \geq 0$, $\pi_{v,u} \in \cusp_{\varrho}(u)$ and $\tau=\pi_{v,u}[s]_{D}$. Let then denote
by $\iota$ the image of $\speh_s(\varrho)$ by the Jacquet-Langlands
correspondence modulo $l$
defined in \cite[\S 1.2.4]{dat-jl}. Then the reduction modulo $l$ of $\tau$ is isomorphic to
\begin{equation} \label{eq-red-tau}
\iota \{-\psfrac{m(\tau)-1}{2}\} \oplus \iota \{-\psfrac{m(\tau)-3}{2}\} \oplus \cdots \oplus \iota \{\psfrac{m(\tau)-1}{2}\}
\end{equation}
where $\iota \{n\}:=\iota \otimes q^{-n \val \circ \nrd}$.

We now recall the notion of level of non degeneracy from \cite[\S4]{zelevinski1}.
The mirabolic subgroup
$M_d(K)$ of $\GL_d(K)$ is the sub-group of matrices with last row $(0,\dots,0,1)$: we denote by
\[
V_d(K)=\{(m_{i,j} \in P_d(K)\sep m_{i,j}= \delta_{i,j} \hbox{ for } j < n\}.
\]
its unipotent radical. We fix a nontrivial character $\psi$ of $K$ and let $\theta$ be the
character of $V_d(K)$ defined by $\theta((m_{i,j}))=\psi(m_{d-1,d})$.
For $G=\GL_r(K)$ or $M_r(K)$, we denote by $\alg(G)$ the abelian category of algebraic
representations of $G$ and, following \cite{zelevinski1}, we introduce
\[
\Psi^-: \alg(M_d(K)) \to \alg(\GL_{d-1}(K), \qquad \Phi^-: \alg (M_d) \to
\alg (M_{d-1}(K))
\]
defined by $\Psi^-=r_{V_d,1}$ (\resp $\Phi^-=r_{V_d,\theta}$) the functor of $V_{d-1}$
coinvariants (resp.\ $(V_{d-1},\theta)$-coinvariants), \cf \cite[1.8]{zelevinski1}.
We also introduce the normalized compact induced functor
\begin{align*}
\Psi^+&:=i_{V,1}: \alg(\GL_{d-1}(K)) \to \alg (M_d(K)),
\\
\Phi^+&:=i_{V,\theta}: \alg(M_{d-1}(K)) \to \alg(M_d(K)).
\end{align*}

\skpt
\begin{prop}[{\cite[p.\,451]{zelevinski1}}]
\begin{itemize}
\item The functors $\Psi^-$, $\Psi^+$, $\Phi^-$ and $\Phi^+$ are exact.

\item $\Phi^- \circ \Psi^+=\Psi^- \circ \Phi^+=0$.

\item $\Psi^-$ (\resp $\Phi^+$) is left adjoint to $\Psi^+$ (\resp $\Phi^-$) and the
following adjunction maps
\[
\Id \to \Phi^- \Phi^+, \qquad \Psi^+ \Psi^- \to \Id,
\]
are isomorphisms, with an exact sequence
\[
0 \to \Phi^+ \Phi^- \to \Id \to \Psi^+ \Psi^-
\to 0.
\]
\end{itemize}
\end{prop}

\begin{defin}
For $\tau \in \alg(M_d(K))$, the representation
\[
\tau^{(k)}:=\Psi^- \circ (\Phi^-)^{k-1}(\tau)
\]
is called the $k$-th derivative of $\tau$. If $\tau^{(k)}\neq 0$ and $\tau^{(m)}=0$
for all $m > k$, then $\tau^{(k)}$ is called the highest derivative of $\tau$.
\end{defin}

\begin{nota}[{\cf \cite[4.3]{zelevinski2}}] \label{nota-nondegeneracy}
Let $\pi \in \alg(\GL_d(K))$ (or $\pi \in \alg(M_d(K)$). The maximal number $k$ such that
$(\pi_{|M_d(K)})^{(k)} \neq (0)$ is called the level of non-degeneracy of $\pi$ and
denoted by $\lambda(\pi)$. We can also iterate the construction so that at the end we obtain
a partition $\underline{\lambda(\pi)}$ of $d$.
\end{nota}

\begin{defin}
A representation $\pi$ of $\GL_d(K)$, over $\overline \Qm_l$ or $\overline \Fm_l$, is then said generic
if its level of non degeneracy $\lambda(\pi)$ is equal to $d$.
\end{defin}

\begin{remark}
Over $\overline \Qm_l$, an irreducible generic representation of
$\GL_d(K)$ looks like $\st_{t_1}(\pi_1) \times \cdots \times \st_{t_r}(\pi_r)$,
where $\pi_1,\dots, \pi_r$ are irreducible cuspidal representations. Note
moreover that the reduction modulo $l$ of any irreducible generic representation
admits a unique generic irreducible subquotient.
\end{remark}

In the following we will be interested in representations
$\st_{t_1}(\chi_1) \times \cdots \times \st_{t_r}(\chi_r)$ where $\chi_1,\dots,
\chi_r$ are unramified characters.

\begin{nota} \label{nota-iwh0}
Associated to a partition $\underline d=(d_1 \geq d_2 \geq d_s)$ of
$d=\sum_{i=1}^s d_i$, we~consider the following Iwahori type subgroup of
$\GL_d(\OC_K)$:
\[
\Iw_v(\underline d):=
\bigl \{g \in \GL_d(\OC_K) \sep (g \bmod \varpi) \in
P_{d_1 < d_1+d_2 < \cdots < d}(\kappa) \bigr\}.
\]
\end{nota}

Recall the well-known following lemma.

\begin{lemma} \label{lem-iwahori}
With the previous notations,
$\st_{t_1}(\chi_1) \times \cdots \times \st_{t_r}(\chi_r)$ has nontrivial
invariant vectors under $\Iw(\underline d)$ if, and only if, $\underline d$
is smaller, for the Bruhat order, to the dual partition associated to
$(t_1,\dots,t_r)$.
\end{lemma}

\begin{remark}
Recall that one way to obtain the dual partition is to use Fejer's diagrams.
To a partition $(d_1 \geq d_2 \geq d_r)$ one can associate a Fejer diagram
with rows of respective size $d_1,\dots,d_r$. Then one can read this Fejer diagram
through its columns whose size gives the dual partition associated to
$(d_1 \geq \cdots \geq d_r)$.
\end{remark}

\skpt
\begin{examples}
\begin{itemize}
\item If $t_1=\cdots=t_r=1$ with $r=d$, then the dual partition of
$(1,\dots,1)$ is $(d)$ and $\chi_1 \times \cdots \times \chi_r$ has nontrivial
invariant vectors under $\Iw(d)=\GL_d(\OC_K)$ and so under all the
Iwahori type subgroup $\Iw(\underline d)$.

\item At the opposite $\st_d(\chi)$ has nontrivial invariant vectors under
$\Iw(1,\dots,1)$ and it is the only Iwahori type subgroup with this property.

\item $\speh_d(\chi)$ has nontrivial invariant vectors under $\GL_d(\OC_K)$.

\item $LT_{\chi}(t-1,d-t)$ which can be seen as a subspace of
$\st_t(\chi_v \{\psfrac{t-d}{2}\} \times \speh_{d-t}(\chi_v \{\sfrac{t}{2})$
has nontrivial invariant vectors under $\Iw(\underline d)$, where
$\underline d$ is the dual partition of $(t,1,\dots,1)$, \ie $\underline d=(d-t+1,1,\dots,1)$. Moreover it has no nontrivial invariant
vectors under $\Iw(\underline d')$ for any $\underline d'$ strictly greater
than $\underline d$.
\end{itemize}
\end{examples}

\subsection{Shimura varieties of KHT type}

\label{para-geo}

Let $F=F^+ E$ be a CM field where $E/\Qm$ is quadratic imaginary and
$F^+/\Qm$ is
totally real with a fixed real embedding \hbox{$\tau:F^+ \hto \Rm$}. For a place $v$ of $F$,
we will denote by
\begin{itemize}
\item $F_v$ the completion of $F$ at $v$,

\item $\OC_v$ the ring of integers of $F_v$,

\item $\varpi_v$ a uniformizer,

\item $q_v$ the cardinal of the residue field $\kappa(v)=\OC_v/(\varpi_v)$.
\end{itemize}
Let $B$ be a division algebra with center $F$, of dimension $d^2$ such that at every place~$x$ of~$F$,
either $B_x$ is split or a local division algebra and suppose $B$ provided with an involution of
second kind $*$ such that $*_{|F}$ is the complex conjugation. For any
$\beta \in B^{*=-1}$, denote by $\sharp_\beta$ the involution $x \mto x^{\sharp_\beta}=\beta x^*
\beta^{-1}$ and let $G/\Qm$ be the group of similitudes, denoted by
$G_\tau$ in \cite{h-t}, defined for every $\Qm$-algebra $R$ by
\[
G(R) \simeq \{(\lambda,g) \in R^\times \times (B^{\op} \otimes_\Qm R)^\times \sep
gg^{\sharp_\beta}=\lambda\}
\]
with $B^{\op}=B \otimes_{F,c} F$.
If $x$ is a place of $\Qm$, split $x=yy^c$ in $E$, then
\begin{equation} \label{eq-facteur-v}
G(\Qm_x) \simeq (B_y^{\op})^\times \times \Qm_x^\times \simeq \Qm_x^\times \times
\prod_{z_i} (B_{z_i}^{\op})^\times,
\end{equation}
where, identifying places of $F^+$ over $x$ with places of $F$ over $y$,
$x=\prod_i z_i$ in $F^+$.

\begin{convention}
For a place $x=yy^c$ of $\Qm$ split in $E$ and $z$
a place of $F$ over $y$, we~shall make throughout the text the following abuse of notation by denoting
$G(F_z)$ in place of the factor $(B_z^{\op})^\times$ in the formula \eqref{eq-facteur-v}.
\end{convention}

In \cite{h-t}, the authors justify the existence of some $G$ like above such that moreover
\begin{itemize}
\item if $x$ is a place of $\Qm$ non split in $E$ then $G(\Qm_x)$ is quasi split;

\item the invariants of $G(\Rm)$ are $(1,d-1)$ for the embedding $\tau$ and $(0,d)$ for the others.
\end{itemize}

As in \cite[bottom of p.\,90]{h-t}, a compact open subgroup $U$ of $G(\Am^\oo)$ is said to be
\emph{small enough}
if there exists a place $x$ such that the projection from $U^v$ to $G(\Qm_x)$ does not contain any
element of finite order except identity.

\begin{nota}
Denote by $\IC$ the set of open compact subgroups small enough of $G(\Am^\oo)$.
For $I \in \IC$, write $\sh_{I,\eta} \to \spec F$ for the associated
Shimura variety of Kottwitz-Harris-Taylor type.
\end{nota}

\begin{defin} \label{defi-spl}
Denote by $\spl$ the set of places $v$ of $F$ such that $p_v:=v_{|\Qm} \neq l$ is split in $E$ and let
$B_v^\times \simeq \GL_d(F_v)$. For each $I \in \IC$, we write
$\spl(I)$ for the subset of $\spl$ of places which does not divide $I$.
\end{defin}

In the sequel, $v$ and $w$ will denote places of $F$ in $\spl$.
For such a place $v$,
the scheme $\sh_{I,\eta}$ has a projective model $\sh_{I,v}$ over $\spec \OC_v$
with special fiber $\sh_{I,s_v}$. For $I$ going through $\IC$, the projective system $(\sh_{I,v})_{I\in \IC}$
is naturally equipped with an action of $G(\Am^\oo) \times \Zm$ such that any
$w_v$ in the Weil group $W_v$ of $F_v$ acts by $-\deg (w_v) \in \Zm$,
where $\deg=\val \circ \art^{-1}$ and $\art^{-1}:W_v^{\ab} \simeq F_v^\times$ is the isomorphism of Artin
sending the geometric Frobenius to a uniformizer.

\begin{notas}
For $I \in \IC$, the Newton stratification of the geometric special fiber $\sh_{I,\bar s_v}$ is denoted by
\[
\sh_{I,\bar s_v}=:\sh^{\geq 1}_{I,\bar s_v} \supset \sh^{\geq 2}_{I,\bar s_v} \supset \cdots \supset
\sh^{\geq d}_{I,\bar s_v},
\]
where $\sh^{=h}_{I,\bar s_v}:=\sh^{\geq h}_{I,\bar s_v} - \sh^{\geq h+1}_{I,\bar s_v}$ is an affine
scheme, smooth of pure dimension $d-h$ built up by the geometric
points whose connected part of their Barsotti-Tate group is of rank $h$.
For each $1 \leq h <d$, write
\[
i_{h}:\sh^{\geq h}_{I,\bar s_v} \hto \sh^{\geq 1}_{I,\bar s_v}, \quad
j^{\geq h}: \sh^{=h}_{I,\bar s_v} \hto \sh^{\geq h}_{I,\bar s_v},
\]
and $j^{=h}=i_h \circ j^{\geq h}$.
\end{notas}

Let $\sigma_0:E \hto
\overline{\Qm}_l$ be a fixed embedding and write $\Phi$ for the set of embeddings
$\sigma:F \hto \overline \Qm_l$ whose restriction to $E$ equals $\sigma_0$.
There exists then, \cf \cite[p.\,97]{h-t}, an explicit bijection between irreducible algebraic representations
$\xi$ of $G$ over $\overline \Qm_l$ and $(d+1)$-uples
$\bigl (a_0, (\overrightarrow{a_\sigma})_{\sigma \in \Phi} \bigr)$,
where $a_0 \in \Zm$ and for all $\sigma \in \Phi$, we have \hbox{$\overrightarrow{a_\sigma}=
(a_{\sigma,1} \leq \cdots \leq a_{\sigma,d})$}.
We~then denote by
\[
V_{\xi,\overline \Zm_l}
\]
the associated $\overline \Zm_l$-local system on $\sh_\IC$.
Recall that an irreducible automorphic representation $\Pi$ is said $\xi$-cohomological if there exists
an integer $i$ such that
\[
H^i \bigl ((\lie ~G(\Rm)) \otimes_\Rm \Cm,U,\Pi_\oo \otimes \xi^\vee \bigr) \neq (0),
\]
where $U$ is a maximal open compact subgroup modulo the center of $G(\Rm)$.
Let $d_\xi^i(\Pi_\oo)$ be the dimension of this cohomology group.

\Subsection{Cohomology of the Newton strata}
\label{para-hecke}

\begin{nota} \label{nota-hixi}
For $1 \leq h \leq d$, let $\IC_v(h)$ be the set of open compact subgroups
\[
U_v(\underline m,h):=
U_v(\underline m^v) \times \left (\begin{array}{cc} I_h & 0 \\ 0 & K_v(m_1) \end{array} \right),
\]
where
\[
K_v(m_1)=\ker \bigl (\GL_{d-h}(\OC_v) \to \GL_{d-h}(\OC_v/ (\varpi_v^{m_1})) \bigr).
\]
We then denote by $[H^i(h,\xi)]$ (\resp $[H^i_!(h,\xi)]$) the image of
\[
\varinjlim_{{I \in \IC_v(h)}} H^i(\sh_{I,\bar s_v,1}^{\geq h}, V_{\xi,\overline \Qm_l}[d-h]),
\qquad \hbox{\resp }
\varinjlim_{{I \in \IC_v(h)}} H^i(\sh_{I,\bar s_v,1}^{\geq h}, j^{\geq h}_{1,!} V_{\xi,\overline \Qm_l}[d-h])
\]
inside the Grothendieck $\groth(v,h)$ of admissible representations of
\[
G(\Am^{\oo}) \times \GL_{d-h}(F_v) \times \Zm.
\]
\end{nota}

\begin{remark}
An element $\sigma \in W_v$ acts through $-\deg \sigma \in \Zm$ and $\Pi_{p_v,0}(\art^{-1} (\sigma))$.
We~moreover consider the action of $\GL_{h}(F_v)$ through
$\val \circ \det: \GL_{h}(F_v) \to \Zm$ and finally
$P_{h,d}(F_v)$ through its Levi factor $\GL_{h}(F_v) \times \GL_{d-h}(F_v)$, \ie its unipotent radical acts trivially.
\end{remark}

From \cite[Prop.\,3.6]{boyer-imj}, for any irreducible tempered automorphic representation
$\Pi$ of $G(\Am)$ and for every $i \neq 0$, the $\Pi^{\oo,v}$-isotypic component of
$[H^i(h,\xi)]$ and $[H^i_!(h,\xi)]$ are zero. About the case $i=0$, for $\Pi$ an irreducible
automorphic tempered $\xi$\nobreakdash-coho\-mo\-logical representation, its local component at $v$ is generic and so looks like
\[
\Pi_v \simeq \st_{t_1}(\pi_{v,1}) \times \cdots \times \st_{t_u}(\pi_{v,u}),
\]
where for $i=1,\dots,u$, $\pi_{v,i}$ is an irreducible cuspidal representation of $\GL_{g_i}(F_v)$.

\begin{prop}[{\cf \cite[Prop.\,3.9]{boyer-imj}}]\label{prop-temperee-explicite}
With the previous notations, we order the representations~$\pi_{v,i}$ such that the first $r$ ones are unramified
characters. Then the $\Pi^{\oo,v}$-isotypic component of $[H^0(h,\xi)]$ is equals to
\[
\Bigl (\frac{\sharp \ker^1(\Qm,G)}{d} \sum_{\Pi' \in \UC_G(\Pi^{\oo,v})}
m(\Pi') d_\xi(\Pi'_\oo) \Bigr) \Bigl (\sum_{1 \leq k \leq r:~ t_k=h} \Pi_v^{(k)} \otimes \chi_{v,k}
\chi^{\psfrac{d-h}{2}} \Bigr),
\]
where
\begin{itemize}
\item $ \ker^1(\Qm,G)$ is the subset of elements of $H^1(\Qm,G)$ which become trivial
in $H^1(\Qm_{p'},G)$ for every prime $p'$;

\item $\Pi_v^{(k)}:= \st_{t_1}(\chi_{v,1}) \times \cdots \times \st_{t_{k-1}}(\chi_{v,k-1}) \times
\st_{t_{k+1}}(\chi_{v,k+1}) \times \cdots \times \st_{t_u}(\chi_{v,u})$ and

\item $\Xi:\frac{1}{2} \Zm \to \overline \Zm_l^\times$
is defined by $\Xi(\sfrac{1}{2})=q_v^{\sfrac{1}{2}}$.

\item $\UC_G(\Pi^{\oo,v})$ is the set of equivalence classes of irreducible automorphic
representations $\Pi'$ of $G(\Am)$ such that $(\Pi')^{\oo,v} \simeq \Pi^{\oo,v}$.
\end{itemize}
\end{prop}

\begin{remark}
With the notations of \cite{boyer-imj}, we only consider Harris-Taylor perverse
sheaves associated to the trivial representation which then selects the unramified
characters to describe its cohomology groups.

Note also that if $[H^0(h,\xi)]$ has nontrivial invariant vectors under some open
compact subgroup $I \in \IC_v(h)$ which is maximal at $v$, then the local component
of $\Pi$ at $v$ is of the following shape $\st_h(\chi_{v,1}) \times \chi_{v,2} \times \cdots \chi_{v,d-h}$,
where the $\chi_{v,i}$ are unramified characters.
\end{remark}

\begin{defin}
For a finite set $S$ of places of $\Qm$ containing the places where $G$ is ramified,
denote by $\Tm^S_{\abs}:=\prod_{x \not \in S} \Tm_{x,\abs}$ the
abstract unramified Hecke algebra, where
$\Tm_{x,\abs} \simeq \overline \Zm_l[X^{\un}(T_x)]^{W_x}$ for $T_x$ a split torus,
$W_x$ is the spherical Weyl group and $X^{\un}(T_x)$ is the set of $\overline \Zm_l$-unramified
characters of $T_x$.
\end{defin}

\begin{example}
For $w \in \spl$, we have
\[
\Tm_{w,\abs}=\overline \Zm_l \bigl [T_{w,i}:~ i=1,\dots,d \bigr ],
\]
where $T_{w,i}$ is the characteristic function of
\[
\GL_d(\OC_w) \diag(\overbrace{\varpi_w,\dots,\varpi_w}^{i}, \overbrace{1,\dots,1}^{d-i})
\GL_d(\OC_w) \subset \GL_d(F_w).
\]
\end{example}

\begin{nota}
Let $\Tm^S_\xi$ be the image of $\Tm^S_{\abs}$ inside
\[
\bigoplus_{i =0}^{2d-2}
\varinjlim_{{I}} H^i(\sh_{I,\bar \eta},V_{\xi,\overline \Qm_l}),
\]
where the limit is taken over
the ideals $I$ which are maximal at each place outside~$S$. For an open compact subgroup $I$ maximal at each place outside $S$,
we will also denote by $\Tm^S_{I,\xi}$ the image of $\Tm^S_{\abs}$ inside
$H^{d-1}(\sh_{I,\bar \eta},V_{\xi,\overline \Qm_l})$.
\end{nota}

Let state some remarks about these Hecke algebras.
\begin{itemize}
\item The torsion cohomology classes give also sets of Satake parameters
and we can ask if they correspond to maximal ideals of $\Tm_\xi^S$ which
we choose to define through the $\overline \Qm_l$\nobreakdash-coefficients.
In \cite{boyer-imj}, we proved that the torsion cohomology classes of
$H^i(\sh_{I,\bar \eta},V_{\xi,\overline \Zm_l})$ with
$I$ maximal at each place outside $S$, can be lifted in characteristic zero.
More precisely, for any place $v' \in \spl$ not in $S$, there exists
a maximal ideal~$\mathfrak m$ of
$(\Tm^{S \cup \{v'\} }_{I\Iw_{v'},\xi})_{\mathfrak m}$,
such that the set of Satake parameters of our torsion cohomology class
is associated to $\mathfrak m$, where $\Iw_{v'}$ is the Iwahori subgroup of $\GL_d(\OC_{v'})$.

\item As explained in the introduction, we will consider maximal ideals
$\mathfrak m$ such that~$\overline \rho_{\mathfrak m}$ is irreducible, so that
the $\overline \Qm_l$-cohomology groups are all concentrated in middle degree,
\ie in degree $0$ if we deal with perverse sheaves.

\item With the notations of Section \ref{para-HT} about the Harris-Taylor local systems,
in \cite{boyer-compositio}, we proved that, except for the $\GL_d(F_v)$-action,
the irreducible subquotients of the
$\overline \Qm_l$-cohomology groups of $j^{\geq tg}_! HT(\pi_v,\Pi_t)$ or
$\lexp p j^{\geq tg}_{!*} HT(\pi_v,\Pi_t)$ are also subquotients of
the cohomology of $\sh_{I,\bar \eta}$.
\end{itemize}

The minimal prime ideals of $\Tm^S_{\xi}$ are the prime ideals above the zero ideal of
$\overline \Zm_l$ and are then in bijection with the prime ideals of
$\Tm^S_{\xi} \otimes_{\overline \Zm_l} \overline \Qm_l$.
To such an ideal, which corresponds
to giving a collection of Satake parameters, is then associated a unique near equivalence class in
the sense of \cite{y-t}, denoted by $\Pi_{\widetilde{\mathfrak m}}$, which is the finite set of
irreducible automorphic cohomological representations whose multi-set of Satake parameters
at each place $x \in \unr(I)$, is given by the multi-set $S_{\widetilde{\mathfrak m}}(x)$ of roots of
the Hecke polynomial
\[
P_{\widetilde{\mathfrak{m}},w}(X):=\sum_{i=0}^d(-1)^i q_w^{\sfrac{i(i-1)}{2}}\,
T_{w,i,\widetilde{\mathfrak m}} X^{d-i} \in \overline \Qm_l[X],
\]
\ie\[
S_{\widetilde{\mathfrak{m}}}(w) := \bigl \{\lambda \in \Tm^S_{\xi} \otimes_{\overline \Zm_l}
\overline \Qm_l / \widetilde{\mathfrak m} \simeq \overline \Qm_l \sep
P_{\widetilde{\mathfrak{m}},w}(\lambda)=0 \bigr\}.
\]
Thanks to \cite{h-t} and \cite{y-t}, we denote by
\[
\rho_{\widetilde{\mathfrak m}}:\gal(\overline F/F) \to \GL_d(\overline \Qm_l)
\]
the Galois representation associated to any $\Pi \in \Pi_{\widetilde{\mathfrak m}}$.
Recall that the reduction modulo~$l$ of $\rho_{\widetilde{\mathfrak m}}$ depends only
of $\mathfrak m$, and was denoted above by $\overline{\rho_{\mathfrak m}}$.
For every $w \in \spl(I)$, we~also denote by $S_{\mathfrak{m}}(w)$ the multi-set of Satake parameters modulo $l$ at $w$ given as the multi-set of roots of
\[
P_{\mathfrak{m},w}(X):=\sum_{i=0}^d(-1)^i q_w^{\sfrac{i(i-1)}{2}}\, \overline{T_{w,i}} X^{d-i}
\in \overline \Fm_l[X],
\]
\ie\[
S_{\mathfrak{m}}(w) := \bigl \{\lambda \in \Tm^S_\xi/\mathfrak m \simeq \overline \Fm_l
\sep P_{\mathfrak{m},w}(\lambda)=0 \bigr\}.
\]

\section{About the nearby cycles perverse sheaf}

Our strategy to compute the cohomology of the KHT-Shimura variety $\sh_{I,\bar \eta}$ with
coefficients in $V_{\xi,\overline \Zm_l}$, is to realize it as the outcome of the nearby cycles
spectral sequence at some place $v \in \spl$.

Note that the role of the local system $V_{\xi,\overline \Zm_l}$
associated to $\xi$ is completely harmless when
dealing with sheaves: one just has to add a tensor product with it to all the
statements without the index $\xi$. In the following we will sometimes not mention the index
$\xi$ in the statements to make formulas more readable. Of course, when looking at the
cohomology groups, the role of $V_{\xi,\overline \Zm_l}$ is crucial as it selects the
automorphic representations which contribute to the cohomology.

\subsection{The case where the level at $v$ is maximal}
\label{para-max}

By the smooth base change theorem, we have
$H^i(\sh_{I,\bar \eta_v}, V_\xi) \simeq H^i(\sh_{I,\bar s_v},V_\xi)$.
As for each $h$ such that $1 \leq h \leq d-1$, the open Newton stratum $\sh^{=h}_{I,\bar s_v}$ is affine,
then $H^i_c(\sh^{=h}_{I,\bar s_v},V_{\xi,\overline \Zm_l}[d-h])$ is zero for $i<0$ and free for $i=0$.
Using this property and the following short exact sequence
of free perverse sheaves
\begin{multline*}
0 \to i_{h+1 \to h,*} V_{\xi,\overline \Zm_l,|\sh_{I,\bar s_v}^{\geq h+1}}[d-h-1] \to
j^{\geq h}_! j^{\geq h,*} V_{\xi,\overline \Zm_l,|\sh^{\geq h}_{I,\bar s_v}}[d-h] \\ \to
V_{\xi,\overline \Zm_l,|\sh^{\geq h}_{I,\bar s_v}}[d-h] \to 0,
\end{multline*}
where $i_{h+1 \to h}: \sh_{I,\bar s_v}^{\geq h+1} \hto
\sh_{I,\bar s_v}^{\geq h}$,
we then obtain for every $i>0$ an exact sequence
\begin{equation} \label{eq-sec}
0 \to H^{-i-1}(\sh^{\geq h}_{I,\bar s_v},V_{\xi,\overline \Zm_l}[d-h]) \to
H^{-i}(\sh^{\geq h+1}_{I,\bar s_v},V_{\xi,\overline \Zm_l}[d-h-1]) \to 0,
\end{equation}
and for $i=0$, a long exact sequence
\begin{multline} \label{eq-sec0}
0 \to H^{-1}(\sh^{\geq h}_{I,\bar s_v},V_{\xi,\overline \Zm_l}[d-h]) \to
H^{0}(\sh^{\geq h+1}_{I,\bar s_v},V_{\xi,\overline \Zm_l}[d-h-1])\\
\to H^0(\sh^{\geq h}_{I,\bar s_v},j^{\geq h}_! j^{\geq h,*} V_{\xi,\overline \Zm_l}[d-h]) \to
H^{0}(\sh^{\geq h}_{I,\bar s_v},V_{\xi,\overline \Zm_l}[d-h]) \to \cdots
\end{multline}
In \cite{boyer-imj}, arguing by induction from $h=d$ to $h=1$, we prove that for a maximal
ideal~$\mathfrak m$ of $\Tm^S_{\xi}$ such that $S_{\mathfrak m}(v)$ does not contain any subset
of the form $\{\alpha,q_v \alpha\}$, all~the cohomology groups
$H^i(\sh_{I,\bar s_v}^{\geq h},V_{\xi,\overline \Zm_l})_{\mathfrak m}$ are free: note that in order
to deal with $i\geq 0$, one has to use Grothendieck-Verdier duality.

Without this hypothesis, arguing similarly, we conclude that any torsion cohomology class
comes from a non strict map
\begin{equation} \label{eq-map-strict-fini}
H^{0}_{\free}(\sh^{\geq h+1}_{I,\bar s_v},V_{\xi,\overline \Zm_l}[d-h-1])_{\mathfrak m} \to
H^0(\sh^{\geq h}_{I,\bar s_v},j^{\geq h}_! j^{\geq h,*} V_{\xi,\overline \Zm_l}[d-h])_{\mathfrak m}.
\end{equation}
In particular it lifts in characteristic zero to some free subquotient of
\[
H^0(\sh^{\geq h}_{I,\bar s_v},j^{\geq h}_! j^{\geq h,*} V_{\xi,\overline \Zm_l}[d-h])_{\mathfrak m}.
\]

\begin{mainassumption}
We argue by contradiction and we suppose there exists a finite level~$I$ maximal at $v$
such that there
exist nontrivial torsion cohomology classes in the $\mathfrak m$-localized cohomology of
$\sh_{I,\bar \eta_v}$ with coefficients in $V_{\xi,\overline \Zm_l}$.
We then fix such a finite level $I$.
\end{mainassumption}

\begin{prop}[{\cf \cite[Lem.\,4.13]{boyer-imj}}] \label{prop-h0I}
Consider $h_0(I)$ maximal such that there exists $i \in \Zm$ with
$H^{d-h_0(I)+i}(\sh_{I,\bar s_v}^{\geq h_0(I)}, V_{\xi,\overline \Zm_l})_{\mathfrak m,\tor}
\neq (0)$. Then we have the following properties:
\begin{itemize}
\item $i=0,1$;

\item for all $1 \leq h \leq h_0(I)$ and $i<h-h_0(I)$,
\[
H^{d-h+i}(\sh_{I,\bar s_v}^{\geq h} V_{\xi,\overline \Zm_l})_{\mathfrak m,\tor} =(0),
\]
while for $i=h-h_0(I)$ it is nontrivial.
\end{itemize}
\end{prop}

\begin{remark}
Note that any system of Hecke eigenvalues
$\mathfrak m$ of $\Tm^S_\xi$ inside the torsion of some
$H^i(\sh_{I,\bar \eta_v},V_{\xi,\overline \Zm_l})$ lifts in characteristic zero, \ie is associated to a minimal prime ideal $\widetilde{\mathfrak m}$ of
$\Tm^{S \cup \{v\}}_{\xi}$: the level of $\widetilde{\mathfrak m}$ outside $v$ can still
be taken to be $I^v$, and, at $v$, using the remark following Proposition
\ref{prop-temperee-explicite}, $\pi_{\widetilde{\mathfrak m},v} \simeq
\st_{h_0(I)+1}(\chi_v) \times \chi_{v,1} \times \cdots \times \chi_{v,d-h_0(I)-1}$,
where $\chi_v, \chi_{v,1},\dots, \chi_{v,d-h_0(I)-1}$ are characters of $F_v^\times$.
\end{remark}

\subsection{Harris-Taylor perverse sheaves over $\overline \Zm_l$}
\label{para-HT}

Consider now the ideals $I^v(n):=I^vK_v(n)$, where
$K_v(n):=\ker(\GL_d(\OC_v) \onto \GL_d(\OC_v/\MC_v^n))$.
Recall then that $\sh_{I^v(n),\bar s_v}^{=h}$ is geometrically induced
under the action of the parabolic subgroup $P_{h,d}(\OC_v/\MC_v^n)$, defined as the
stabilizer of the first $h$ vectors of the canonical basis of $F_v^d$. Concretely this means there
exists a closed subscheme $\sh_{I^v(n),\bar s_v,\overline{1_{h}}}^{=h}$ stabilized by the Hecke
action of $P_{h,d}(F_v)$ and such that
\[
\sh_{I^v(n),\bar s_v}^{=h} = \sh_{I^v(n),\bar s_v,\overline{1_{h}}}^{=h}
\times_{P_{h,d}(\OC_v/\MC_v^n)} \GL_d(\OC_v/\MC_v^n),
\]
meaning that $\sh_{I^v(n),\bar s_v}^{=h} $ is the disjoint union of copies of
$\sh_{I^v(n),\bar s_v,\overline{1_{h}}}^{=h}$ indexed by
$\GL_d(\OC_v/\MC_v^n)/P_{h,d}(\OC_v/\MC_v^n)$ and
exchanged by the action of
$\GL_d(\OC_v/\MC_v^n)$.

\begin{nota} \label{nota-strata}
For any $g \in \GL_d(\OC_v/\MC_v^n)/P_{h,d}(\OC_v/\MC_v^n)$, we denote by
$\sh^{=h}_{I^v(n),\bar s_v,g}$ the \emph{pure} Newton stratum defined as the image of
$\sh_{I^v(n),\bar s_v,\overline{1_{h}}}^{=h}$ by $g$.
Its closure in $\sh_{I^v(n),\bar s_v}$ is then denoted by
$\sh^{\geq h}_{I^v(n),\bar s_v,g}$.
\end{nota}

Let then denote by $\mathfrak m^v$ the multiset of
Hecke eigenvalues given by $\mathfrak m$ but outside $v$ and
introduce for any representation $\Pi_h$ of $\GL_h(F_v)$:
\[
H^i(\sh_{I^v(\oo),\bar s_v,\overline{1_{h}}}^{\geq h}, V_{\xi,\overline \Zm_l})_{\mathfrak m^v} \otimes \Pi_h:=
\varinjlim_{{n}} H^i(\sh_{I^v(n),\bar s_v,\overline{1_{h}}}^{\geq h},
V_{\xi,\overline \Zm_l})_{\mathfrak m^v} \otimes \Pi_h,
\]
as a representation of $\GL_h(F_v) \times \GL_{d-h}(F_v)$, where $g \in \GL_h(F_v)$
acts on $\Pi_h$ as well as on $H^i(\sh_{I^v(n),\bar s_v,\overline{1_{h}}}^{\geq h},
V_{\xi,\overline \Zm_l})_{\mathfrak m^v}$ through the determinant map \hbox{$\det: \GL_h(F_v)\onto F_v^\times$}. Note moreover that the unipotent radical of
$P_{h,d}(F_v)$ acts trivially on these cohomology groups.
We then introduce their induced version
\[
H^i(\sh^{\geq h}_{I^v(\oo),\bar s_v},\Pi_h \otimes V_{\xi,\overline \Zm_l})_{\mathfrak m^v}
\simeq \ind_{P_{h,d}(F_v)}^{\GL_d(F_v)} H^i(\sh_{I^v(\oo),\bar s_v,\overline{1_h}}^{\geq h}, V_{\xi,\overline \Zm_l})_{\mathfrak m^v} \otimes \Pi_h.
\]

More generally, with the notations of \cite{boyer-invent2}, replace now the trivial representation by
an irreducible cuspidal representation $\pi_v$ of $\GL_g(F_v)$ for some $1 \leq g \leq d$.

\begin{notas}
Let $1 \leq t \leq s:=\lfloor d/g \rfloor$ and let $\Pi_t$ be any representation of
$\GL_{tg}(F_v)$. We then denote by
\[
\widetilde{HT}_{1,\overline \Zm_l}(\pi_v,\Pi_t):=\LC(\pi_v[t]_D)_{\overline{1_{tg}},\overline \Zm_l}
\otimes \Pi_t \otimes \Xi^{\psfrac{tg-d}{2}}
\]
the $\overline \Zm_l$-Harris-Taylor local system on the Newton stratum
$\sh^{=tg}_{I,\bar s_v,\overline{1_{tg}}}$, where
\begin{itemize}
\item $\LC(\pi_v[t]_D)_{\overline{1_{tg}},\overline \Zm_l}$ is defined in \cite[\S IV-1]{h-t}, by means of
Igusa varieties attached to the representation $\pi_v[t]_D$ of the division
algebra of dimension
$(tg)^2$ over $F_v$ associated to $\st_t(\pi_v)$ by the Jacquet-Langlands correspondence,

\item $\Xi:\frac{1}{2} \Zm \to \overline \Zm_l^\times$ defined by $\Xi(\sfrac{1}{2})=q^{1/2}$.
\end{itemize}
We also introduce the induced version
\[
\widetilde{HT}(\pi_v,\Pi_t)_{\overline \Zm_l}:=
\Bigl (\LC(\pi_v[t]_D)_{\overline{1_{tg}},\overline \Zm_l}
\otimes \Pi_t \otimes \Xi^{\psfrac{tg-d}{2}} \Bigr) \times_{P_{tg,d}(F_v)} \GL_d(F_v),
\]
where the unipotent radical of $P_{tg,d}(F_v)$ acts trivially and the action of
\[
\Bigl(g^{\oo,v},\left (\begin{array}{cc} g_v^c & * \\ 0 & g_v^\mathrm{et} \end{array} \right),\sigma_v\Bigr)
\in G(\Am^{\oo,v}) \times P_{tg,d}(F_v) \times W_v
\]
is given by
\begin{itemize}
\item the action of $g_v^c$ on $\Pi_t$ and
$\deg(\sigma_v) \in \Zm$ on $ \Xi^{\psfrac{tg-d}{2}}$,

\item and the action of $(g^{\oo,v},g_v^\mathrm{et},\val(\det g_v^c)-\deg \sigma_v)
\in G(\Am^{\oo,v}) \times \GL_{d-tg}(F_v) \times \Zm$ on $\LC_{\overline \Qm_l}
(\pi_v[t]_D)_{\overline{1_{tg}},\overline \Zm_l} \otimes \Xi^{\psfrac{tg-d}{2}}$.
\end{itemize}
We also introduce
\[
HT(\pi_v,\Pi_t)_{1,\overline \Zm_l}:=\widetilde{HT}(\pi_v,\Pi_t)_{1,\overline \Zm_l}
[d-tg],
\]
and the perverse sheaf
\[
P(t,\pi_v)_{1,\overline \Zm_l}:=\lexp p j^{=tg}_{\overline{1_{tg}},!*} HT(\pi_v,\st_t(\pi_v))_{1,\overline \Zm_l}
\otimes \Lm(\pi_v),
\]
and their induced version, $HT(\pi_v,\Pi_t)_{\overline \Zm_l}$ and $P(t,\pi_v)_{\overline \Zm_l}$, where
\[
j^{=h}_{\overline{1_h}}=i^h \circ j^{\geq h}_{\overline{1_h}}:\sh^{=h}_{I,\bar s_v,\overline{1_h}} \hto
\sh^{\geq h}_{I,\bar s_v} \hto \sh_{I,\bar s_v}
\]
and $\Lm^\vee$, the dual of $\Lm$, is the local Langlands correspondence
which sends geometric frobenii to uniformizers.
Finally we will also use the index $\xi$ in the notations, for example
$HT_\xi(\pi_v,\Pi_t)_{\overline \Zm_l}$,
when we twist the sheaf with $V_{\xi,\overline \Zm_l}$.
\end{notas}

With the previous notations, from \eqref{eq-red-tau}, we deduce the following equality in the
Grothendieck group of Hecke-equivariant local systems
\begin{equation} \label{eq-chgt-cuspi}
m(\varrho)l^{u}\Bigl [ \Fm \LC_{\xi,\overline \Zm_l}(\pi_{v,u}[t]_D) \Bigr ] =
\Bigl [ \Fm \LC_{\xi,\overline \Zm_l}(\pi_{v,-1}[tm(\varrho)l^u]_D) \Bigr ].
\end{equation}

We now focus on the perverse Harris-Taylor sheaves. Note first, \cf \cite[(2.2--2.4)]{juteau}, that over
$\overline \Zm_l$, there are two notions of intermediate extension associated to
the two classical $t$-structures $p$ and $p+$: essentially they come from
the choice about~$\Fm_l$ as a sheaves over the point, represented by the
complex $\Zm_l \To{\times l} \Zm_l$ and where one decides to
put the zero grading on the second factor, which corresponds to the $p$-structure, or on the first one which gives the $p+$-structure.
So for every $\pi_v \in \cusp_{\varrho}$ of $\GL_g(F_v)$ and $1 \leq t \leq d/g$, we can define:
\begin{equation} \label{eq-bimorphism}
\lexp p j^{=tg}_{!*} HT(\pi_{v,},\Pi_t)_{\overline \Zm_l} \htarrowp \lexp {p+} j^{=tg}_{!*} HT(\pi_{v},\Pi_t)_{\overline \Zm_l},
\end{equation}
the symbol $\htarrowp$ meaning bimorphism, \ie both a monomorphism and an epimorphism, so that
the cokernel for the $t$-structure $p$ (\resp the kernel for $p+$) has support in
$\sh^{\geq tg+1}_{I,\bar s_v}$. When $\pi_v$ is a character, \ie when $g=1$, the associated
bimorphisms are isomorphisms, as explained in the following lemma,
but in general they are not.

\begin{lemma} \label{lem-ext0}
With the previous notations, we have an isomorphism
\[
\lexp p j^{\geq h}_{\overline{1_h},!*} HT(\chi_v,\Pi_h)_{1,\overline \Zm_l}
\simeq \lexp {p+} j^{\geq h}_{\overline{1_h},!*} HT(\chi_v,\Pi_h)_{1,\overline \Zm_l}.
\]
\end{lemma}

\begin{proof}
Recall that $\sh^{\geq h}_{I,\bar s_v,\overline{1_h}}$ is smooth over
$\spec \overline \Fm_p$. Up to a modification of the action of the fundamental
group through the character $\chi_v$, we have
\[
HT(\chi_v,\Pi_h)_{1,\overline \Zm_l}[h-d]= (\overline \Zm_l)_{|\sh^{\geq h}_{I,\bar s_v,\overline{1_h}}} \otimes \Pi_h.
\]
Then $HT(\chi_v,\Pi_h)_{1,\overline \Zm_l}$ is perverse for both $t$-structures, with
\[
i^{h \leq +1,*}_{\overline{1_h}} HT(\chi_v,\Pi_h)_{1,\overline \Zm_l} \in
\lexp p \DC^{< 0} \quad\text{and}\quad
i^{h \leq +1,!}_{\overline{1_h}} HT(\chi_v,\Pi_h)_{1,\overline \Zm_l} \in
\lexp {p+} \DC^{\geq 1}.\qedhere
\]
\end{proof}

\begin{remark}
One of the main results of \cite[Prop.\,2.4.1]{boyer-duke} is the fact that
the previous lemma holds for any $\pi_v \in \cusp_\varrho(-1)$.
As explained in the introduction, with the hypothesis on the order of $q_v$ modulo $l$,
which is supposed to be strictly greater than~$d$, for $\varrho=\mathds 1_v$ being the trivial character,
we do not need to bother about the representations
$\pi_v \in \cusp_{\mathds 1_v}(u)$ for $u \geq 0$,
\cf the remark after Notation \ref{nota-mvarrho}.
\end{remark}

\subsection{Filtrations of the nearby cycles perverse sheaf}
\label{para-fil-psi}

Let us denote by\footnote{We decide not to add $I$ in the list of indexes.}
\[
\Psi_{v}:=R\Psi_{\eta_v}(\overline \Zm_l[d-1])(\psfrac{d-1}{2})
\]
the nearby cycles autodual free perverse sheaf on the geometric special fiber $\sh_{I,\bar s_v}$
of~$\sh_I$. We also set $\Psi_{\xi,v}:=\Psi_{v} \otimes V_{\xi,\overline \Zm_l}$.
Let $\scusp_{\overline \Fm_l}(g)$ be the set of inertial equivalence
classes of irreducible $\overline \Fm_l$-supercuspidal representations of
$\GL_g(F_v)$. In \cite[Prop.\,3.3.4]{boyer-local-ihara}, we prove the following splitting:
\begin{equation} \label{eq-psi-dec}
\Psi_{v} \simeq \bigoplus_{g=1}^d
\bigoplus_{\varrho \in \scusp_{\overline \Fm_l}(g)} \Psi_{\varrho},
\end{equation}
with the property that the irreducible subquotients of
\[
\Psi_{\varrho} \otimes_{\overline \Zm_l} \overline \Qm_l \simeq
\bigoplus_{\pi_v \in \cusp_\varrho} \Psi_{\pi_v}
\]
are exactly the perverse Harris-Taylor sheaves, of level $I$,
associated to an irreducible cuspidal $\overline \Qm_l$-representation
of some $\GL_g(F_v)$ such that the supercuspidal support of the reduction modulo $l$ of
$\pi_v$ is a segment associated to the inertial class $\varrho$.

This splitting relies on the various filtrations defined over $\overline \Zm_l$,
of $\Psi_v$ constructed by means of the Newton stratification, \cf \cite{boyer-FT}.
Using the
adjunction morphism \hbox{$j^{=t}_! j^{=t,*}\to\Id$} as in \cite{boyer-torsion}, we then define
a filtration of $\Psi_{\xi,\varrho}$
\[
\Fil^0_!(\Psi_{\xi,\varrho}) \harrow \Fil^1_!(\Psi_{\xi,\varrho}) \harrow
\Fil^{2}_!(\Psi_{\xi,\varrho})
\harrow \cdots \harrow \Fil^{d}_!(\Psi_{\xi,\varrho}),
\]
where the symbol $\harrow$ means a monomorphism such that the cokernel
is torsion-free, which here means that
$\Fil^{t}_!(\Psi_{\xi,\varrho})$ is the saturated image of
$j^{=t}_!j^{=t,*} \Psi_{\xi,\varrho} \to \Psi_{\xi,\varrho}$.
We~then denote by $\gr^k_!(\Psi_{\xi,\varrho})$ the graded parts and we have a spectral sequence:
\begin{equation} \label{eq-ss2}
E_{!,\varrho,1}^{p,q}=H^{p+q}(\sh_{I,\bar s_v},\gr^{-p}_!(\Psi_{\xi,\varrho})) \implies
H^{p+q}(\sh_{I,\bar s_v},\Psi_{\xi,\varrho}).
\end{equation}

\begin{remark}
Over $\overline \Qm_l$, in \cite{boyer-torsion} we proved that
$\Fil^k_!(\Psi_{\xi,\varrho}) \otimes_{\overline \Zm_l} \overline \Qm_l$
is $\ker N^k$ where $N$ is the monodromy operator at $v$.

We can refine this filtration to define a filtration
$\Fil^{\csbullet} (\gr^k_!(\Psi_{\xi,\varrho}))$ of $\gr^k_!(\Psi_{\xi,\varrho})$ where,
over $\overline \Qm_l$,
$\Fil^{\csbullet} (\gr^k_!(\Psi_{\xi,\pi_v}))$
coincides with the filtration by iterated images of~$N$, \ie $\gr^r(\gr^k_!(\Psi_{\xi,\pi_v}))=\im N^r \cap \ker N^k$, so that
we recover the usual monodromy bi-filtration of \cite{boyer-invent2}.

We then obtain, \cf \cite{boyer-torsion}, an exhaustive filtration
of stratification $\Fill^{\csbullet}(\Psi_{\xi,\varrho})$ of $\Psi_{\xi,\varrho}$ whose graded parts are free,
isomorphic to some free perverse Harris-Taylor sheaves.
Let us denote by $\grr^k(\Psi_{\xi,\varrho}):=\Fill^k(\Psi_{\xi,\varrho})/\Fill^{k-1}(\Psi_{\xi,\varrho})$
the graded parts of this exhaustive filtration.
We then have a spectral sequence
\begin{equation} \label{eq-ss1}
E_1^{p,q}=H^{p+q}(\sh_{I,\bar s_v},\grr^{-p}(\Psi_{\xi,\varrho})) \implies
H^{p+q}(\sh_{I,\bar s_v},\Psi_{\xi,\varrho}),
\end{equation}
where we recall the bimorphisms
\begin{multline} \label{eq-grppsi}
\lexp p j^{=tg}_{!*} HT_\xi(\pi_v,\st_t(\pi_v))_{\overline \Zm_l}(\psfrac{1-t+2i}{2}) \htarrowp
\grr^k(\Psi_{\xi,v})\\
\htarrowp \lexp {p+} j^{=tg}_{!*} HT_\xi(\pi_v,\st_t(\pi_v))_{\overline \Zm_l}(\psfrac{1-t+2i}{2}),
\end{multline}
for some irreducible cuspidal representation $\pi_v$ of $\GL_g(F_v)$ with $1 \leq t \leq d/g$ and
$0 \leq i \leq \lfloor d/g \rfloor -1$ and of type $\varrho$.
\end{remark}

\begin{remark}
In \cite{boyer-duke} we proved that, following the previous process, then all the previous graded parts of $\Psi_{\varrho}$
are isomorphic to $p$-intermediate extensions.
In the following we will only consider the
case where $\varrho$ is the trivial character so that as the order of~$q_v$
being supposed
to be $>d$, then $\cusp_\varrho$ is made of characters
in which case, \cf Lem\-ma~\ref{lem-ext0}, the $p$ and $p+$ intermediate extensions
coincide. Note that in the following
we will not use the results of \cite{boyer-duke}.
\end{remark}

\section{Irreducibility implies freeness}
\label{para-proof}

Recall, \cf the main assumption in Section \ref{para-max},
that we argue by contradiction assuming there exist nontrivial
cohomology classes in some of the
$H^i(\sh_{I,\bar \eta},V_{\xi,\overline \Zm_l})_{\mathfrak m}$.
The strategy is then to choose a place $v \in \spl_I$ such that the order of $q_v$
modulo $l$ is strictly greater than $d$ and to allow ramification at $v$, either infinitely
with $I^v(\oo)$ as denoted in the next paragraph, or
of Iwahori type.

In the arguments we need to consider another place $w \neq v$
verifying the same hypothesis as $v$, \ie $w \in \spl_I$ and such that the order of
$q_w$ modulo $l$ is strictly greater than $d$. We will also allow to
increase infinitely the level at $w$. As the order of both $q_v$ and
$q_w$ is supposed to be strictly greater than $d$, then the functors of invariants
by any open compact subgroups either at $v$ or $w$, are exact. In particular
as there exist nontrivial torsion classes in level $I$ in some degree $i$, when
$I$ is maximal at $v$ and
$w$, there also exist nontrivial torsion classes in degree $i$, whatever the
level~$J$ such that $J^{v,w} = I^{v,w}$ is.
In particular when the level at $v$ is infinite, from the splitting
\begin{multline*}
H^i(\sh_{I^v(\oo),\bar \eta_v},V_{\xi,\overline \Zm_l}[d-1])_{\mathfrak m} \simeq
H^i(\sh_{I^v(\oo),\bar s_v},\Psi_{\xi,v})_{\mathfrak m} \\
\simeq \bigoplus_{g=1}^d
\bigoplus_{\varrho \in \scusp_{\overline \Fm_l}(g)} H^i(\sh_{I^v(\oo),\bar s_v},
\Psi_{\xi,\varrho})_{\mathfrak m},
\end{multline*}
there also exist nontrivial torsion classes
in $H^i(\sh_{I^v(\oo),\bar s_v},\Psi_{\xi,\varrho})_{\mathfrak m}$
when $\varrho=\mathds 1_v$ is the trivial character.

\begin{remark}
From now on, localization at $\mathfrak m$ means that we prescribe
the Satake parameters modulo $l$ as usual, but outside $\{v,w\}$.
\end{remark}

Let now explain the main steps of the following sections.

(a) Following the arguments of the previous section,
we first analyze the torsion cohomology
classes of Harris-Taylor perverse sheaves with infinite level at $v$, and we deduce,
\cf Lemma \ref{lem-rem-after1}, that, as $\overline \Fm_l$-representations
of $\GL_d(F_v)$,
irreducible subquotients of the $l$-torsion
of their cohomology in infinite level at $v$, with highest non degeneracy level, appear in degrees $0,1$.

(b) In Section \ref{para-generic}, considering always infinite level at $v$,
we analyze the torsion
coho\-mology classes of the graded parts $\gr^t_!(\Psi_\varrho)$ of the filtration of
stratification $\Fil^{\csbullet}_!(\Psi_{\varrho})$ and more specifically when
$\varrho=\mathds 1_v$ is the trivial character.
We then deduce, \cf Lem\-ma~\ref{lem-generic-sq}, that the $l$-torsion of
$H^i(\sh_{I^v(\oo),\bar s_v},V_{\xi,\overline \Zm_l})_{\mathfrak m}$
does not have, as an $\overline \Fm_l$\nobreakdash-representation of $\GL_d(F_v)$,
any irreducible generic subquotient whose supercuspidal support is made of
characters.

(c) In Section \ref{para-modification}, we obtain two fundamental results.
\begin{itemize}
\item First, \cf Lemma \ref{lem-important},
under the hypothesis that there exist nontrivial torsion cohomology classes, we show that
the graded pieces $\Gamma_k$ of the filtration given by the spectral sequence of
vanishing cycles, of the free quotient of
$H^0(\sh_{I^v(\oo),\bar \eta_v},V_{\xi,\overline \Zm_l})_{\mathfrak m}$
are not always given by the lattice $\Gamma_0$ of
\[
H^0(\sh_{I,\bar s_v},P_\xi(t,\chi_v)_{\overline \Zm_l}(\psfrac{1-t+2\delta}{2}))_{\mathfrak m}
\otimes_{\overline \Zm_l} \overline \Qm_l
\]
given by the integral cohomology of
$P_\xi(t,\chi_v)_{\overline \Zm_l}$. Roughly, there exist some $k$ and a short exact
sequence $\Gamma_0 \hto \Gamma_k \onto T$, where $T$ is nontrivial and torsion.

\item We then play with the action of $\GL_d(F_w)$ by allowing infinite level at $w$. The main
observation at the end of the section, \cf Proposition \ref{prop-important},
is that, as an $\overline \Fm_l$\nobreakdash-rep\-re\-sentation
of $\GL_d(F_w)$, all the irreducible subquotients of the $l$-torsion of the cokernels~$T$ between two lattices in the previous point, up to multiplicities,
are also subquotients of the $l$-torsion of the cohomology of the Shimura variety.
In particular,
as $v$ and $w$ are playing symmetric roles, these subquotients are not generic,
\cf Corollary \ref{coro-absurd}.
\end{itemize}

(d) In Section \ref{para-final}, the last step is to prove, under the absurd hypothesis that
there exist nontrivial torsion cohomology classes while
$\overline \rho_{\mathfrak m}$ being irreducible, that $S_{\mathfrak m}(v)$
contains a full Zelevinsky line modulo $l$
$\{\lambda q_v^n \in \overline \Fm_l\sep n \in \Zm\}$ which is of order the order
of $q_v$ modulo $l$. As this order is supposed to be strictly greater than $d$,
this is a contradiction. For more insight on the strategy to prove this fact using the previous
properties about lattices, we refer to
the introduction of Section \ref{para-final}.

\subsection{Torsion classes for Harris-Taylor perverse sheaves.}
\label{para-torsion1}

We focus on the torsion in the cohomology groups of the Harris-Taylor perverse
sheaves $P_\xi(\chi_v,t)_{\overline \Zm_l}$ when the level at $v$ is infinite, and as explained above, \cf the main assumption of Section \ref{para-max},
especially when the reduction modulo $l$ of $\chi_v$ is the trivial character.

\begin{nota} \label{nota-mvw}
We will denote by $I^{v} \in \IC$ a finite
level outside $v$, and we also denote by
$\mathfrak m$ the maximal ideal of $\Tm^{S \cup \{v\} }_{\xi}$
associated to $\mathfrak m$, \ie we do not prescribe the Satake
parameters modulo $l$ at $v$.
Let us also set\vspace*{-3pt}
\[
H^i(\sh_{I^v(\oo),\bar \eta_v},V_{\xi,\overline \Zm_l})_{\mathfrak m}:=
\varinjlim_{{I_v}}
H^i(\sh_{I^vI_v,\bar \eta_v},V_{\xi,\overline \Zm_l})_{\mathfrak m},
\]
which can be viewed as a $\overline \Zm_l[\GL_d(F_v)]$-module.
Then, morally, $I^v(\oo)$ is a finite level outside $v$ and infinite at $v$.
\end{nota}

\begin{prop}[\cf \cite{boyer-duke}, second global result of the introduction]
We have
the following resolution of $\lexp p j_{!*}^{=t} HT(\chi_{v},\Pi_t)_{\overline \Zm_l}$\vspace*{-3pt}
\begin{multline} \label{eq-resolution0}
0 \to j_!^{=d} HT(\chi_{v},\Pi_t \{\psfrac{t-s}{2}\}) \times
\speh_{d-t}(\chi_{v}\{t/2\}))_{\overline \Zm_l}
\otimes \Xi^{\psfrac{s-t}{2}} \to \cdots \\
\to j_!^{=t+1} HT(\chi_{v},\Pi_t \{-1/2\} \times \chi_{v} \{t/2\})_{\overline \Zm_l}
\otimes \Xi^{\sfrac{1}{2}} \\
\to j_!^{=t} HT(\chi_{v},\Pi_t)_{\overline \Zm_l}
\to \lexp p j_{!*}^{=t} HT(\chi_{v},\Pi_t)_{\overline \Zm_l} \to 0.
\end{multline}
\end{prop}

Note that
\begin{itemize}
\item as this resolution is equivalent to the computation of the sheaves
cohomology groups
of $\lexp p j_{!*}^{=h} HT(\chi_v,\st_h(\chi_v))_{\overline \Zm_l} $
as explained for example in
\cite[Prop.\,B.1.5]{boyer-duke}, then,
over $\overline \Qm_l$, the proposition follows from the main results of \cite{boyer-invent2}.

\item Over $\overline \Zm_l$, as every terms are free perverse sheaves, then
all maps are necessary strict.

\item This resolution, for a a general supercuspidal representation
with supercuspidal reduction modulo $l$, is one of the main result
of \cite[\S2.3]{boyer-duke}. However, in the case of a character $\chi_v$ the arguments are much easier.
\end{itemize}

Consider the finite level $I^v(n)=I^v I_{v,n}$, where
\[
I_{v,n}=\ker (\GL_d(\OC_v) \onto
\GL_d(\OC_v/\varpi_v^n).
\]
The strata
$\sh^{\geq h}_{I^v(n),\bar s_v,\overline{1_h}}$ are smooth, then, \cf the proof of Lemma
\ref{lem-ext0}, the constant sheaf, up to shift, is perverse and
so equal to the intermediate extension of the constant sheaf, shifted by $d-h$,
on $\sh^{=h}_{I^v(n),\bar s_v,\overline{1_h}}$. In particular, as a constant sheaf,
its sheaf cohomology groups are well-known,
so, over $\sh^{\geq h}_{I^v(n),\bar s_v,\overline{1_h}}$ and so
for $\sh^{\geq h}_{I^v(\oo),\bar s_v,\overline{1_h}}$, the resolution is completely obvious for
$\lexp p j_{\overline{1_h},!*}^{=h} HT(\chi_v,\st_h(\chi_v))_{1,\overline \Zm_l}$
if one remembers that $\speh_i(\chi_v)$ is just the character $\chi_v \circ \det$
of $\GL_i(F_v)$.

The stated resolution is then simply the induced version of the resolution
of $\lexp p j^{=h}_{\overline{1_h},!*} HT(\chi_v,\st_h(\chi_v))_{1,\overline \Zm_l}$:
recall that a direct sum of intermediate extensions is still an intermediate
extension.

By the adjunction property, the map
\begin{multline} \label{eq-map1}
j_!^{=t+\delta} HT(\chi_{v},\Pi_t \{\sfrac{-\delta}{2}\}) \times \speh_{\delta}
(\chi_{v}\{t/2\}))_{\overline \Zm_l} \otimes \Xi^{\delta/2} \\
\to j_!^{=t+\delta-1} HT(\chi_{v},\Pi_t \{\psfrac{1-\delta}{2}\}) \times
\speh_{\delta-1}(\chi_{v}\{t/2\}))_{\overline \Zm_l} \otimes \Xi^{\psfrac{\delta-1}{2}}
\end{multline}
is given by
\begin{multline}\label{eq-map2}
HT(\chi_{v},\Pi_t \{\sfrac{-\delta}{2}\} \times \speh_{\delta}(\chi_{v}\{t/2\}))_{\overline \Zm_l}
\otimes \Xi^{\delta/2} \to \\
\lexp p i^{t+\delta,!} j_!^{=t+\delta-1} HT(\chi_{v},\Pi_t \{\psfrac{1-\delta}{2}\}) \times
\speh_{\delta-1}(\chi_{v}\{t/2\}))_{\overline \Zm_l} \otimes \Xi^{\psfrac{\delta-1}{2}}.
\end{multline}
We then have
\begin{multline} \label{eq-map3}
\lexp p i^{t+\delta,!} j_!^{=t+\delta-1} HT(\chi_{v},\Pi_t \{\psfrac{1-\delta}{2}\}) \times
\speh_{\delta-1}(\chi_{v}\{t/2\}))_{\overline \Zm_l} \otimes \Xi^{\psfrac{\delta-1}{2}} \\ \simeq
HT \Bigl (\chi_{v}, \Pi_t \{\psfrac{1-\delta}{2}\}) \times \bigl (\speh_{\delta-1}(\chi_{v} \{-1/2\})
\times \chi_{v} \{\psfrac{\delta-1}{2}\} \bigr) \{t/2\} \Bigr)_{\overline \Zm_l} \otimes \Xi^{\delta/2}.
\end{multline}
Indeed one can compute $\lexp p i^{h+1,!} j_{!}^{=h} HT(\chi_v,\Pi_h)_{\overline \Zm_l}$ by means of
the spectral sequence associated to the exhaustive filtration of stratification of
$j_{!}^{=h} HT(\chi_v,\Pi_h)_{\overline \Zm_l}$
\begin{multline} \label{eq-filj}
(0)=\Fil^0(\chi_{v},h) \hto \Fil^{-d}(\chi_{v},h) \hto \cdots
\\
\hto \Fil^{-h}(\chi_{v},h)=j^{=h}_{!} HT(\chi_{v},\Pi_h)_{\overline \Zm_l}
\end{multline}
with graded parts, using Lemma \ref{lem-ext0} and \cite{boyer-torsion},
\[
\gr^{-k}(\chi_{v},h) \simeq \lexp p j^{=k}_{!*}
HT(\chi_{v},\Pi_h \{\psfrac{h-k}{2}\} \otimes \st_{k-h}(\chi_{v}\{h/2\}))_{\overline \Zm_l}(\psfrac{h-k}{2}).
\]
As remarked before, the sheaf cohomology groups of
\[
i^{h+1,*} \lexp p j_{!*}^{h+k} HT(\chi_v,\Pi_h \{-k/2\} \times \st_k(\chi_v)(h/2))_{\overline \Zm_l}
\]
are torsion-free, so, by Grothendieck-Verdier duality, the same is true for
\[
i^{h+1,!} \lexp p j_{!*}^{h+k} HT(\chi_v,\Pi_h \{-k/2\}\times \st_k(\chi_v)(h/2))_{\overline \Zm_l}.
\]
The statement follows then from the fact that, over
$\overline \Qm_l$, the previous spectral sequence degenerates at $E_1$.

\begin{remark}
This property is also true when we replace the character $\chi_v$ by
any irreducible cuspidal representation $\pi_v$, \cf \cite{boyer-duke}.
\end{remark}

\begin{fact}
In particular, up to homothety, the map \eqref{eq-map3}, and so those of \eqref{eq-map2}, is
unique. Finally, as the maps of \eqref{eq-resolution0} are strict, the given maps \eqref{eq-map1}
are uniquely determined, that is, if we forget the infinitesimal parts these maps are independent
of the chosen $t$ in \eqref{eq-resolution0}.
\end{fact}

We now copy the arguments of Section \ref{para-max}.

\begin{nota} \label{nota-ih}
For every $h$ such that $1 \leq h \leq d$, let us denote by $i_{I^v}(h,\chi_v)$ the smallest index $i$
such that
$H^i(\sh_{I^{v}(\oo),\bar s_v},\lexp p j^{=h}_{!*} HT_\xi(\chi_{v},\Pi_h)_{\overline \Zm_l})_{\mathfrak m}$
has nontrivial torsion: if it doesn't exists, then set $i_{I^v}(h,\chi_v)=+\oo$.
\end{nota}

\begin{remark}
By duality, as
$\lexp p j_{!*}^{=h}=\lexp {p+} j^{=h}_{!*}$ for Harris-Taylor local systems
associated to a
character, note that when $i_I(h,\chi_v)$ is finite then $i_{I^v}(h,\chi_v) \leq 0$.
\end{remark}

\begin{nota} \label{nota-h0}
Suppose there exists $I \in \IC$ such that
there exists $h$ with $1 \leq h \leq d$ and with $i_{I^v}(h,\chi_v)$ finite, and denote by
$h_0(I^v,\chi_v)$
the biggest such $h$.
\end{nota}

\begin{lemma} \label{lem-torsion-min}
For $1 \leq h \leq h_0(I^v,\chi_v)$ then we have
\[
i_{I^v}(h,\chi_v)=h-h_0(I^v,\chi_v).
\]
Moreover $\frob_v$ acts by $\chi_v(\frob_v) q_v^{\psfrac{h_0(I^v)+1-h}{2}}$.
\end{lemma}

\begin{proof}
Note first that for every $h$ such that $h_0(I^v,\chi_v) \leq h \leq s$, the cohomology groups of
$j^{=h}_! HT_\xi(\chi_{v},\Pi_h)$ are torsion-free.
The spectral sequence associated to the filtration
\eqref{eq-filj}, localized at $\mathfrak m$, is then concentrated in
middle degree and is torsion-free.

Consider then the spectral sequence associated to the resolution \eqref{eq-resolution0}: its
$E_1$ terms are torsion-free and it degenerates at $E_2$.
As, by hypothesis, the abutment of this spectral sequence is free
and is equal to only one $E_2$ terms, we deduce that all the maps
\begin{multline} \label{eq-map1-coho0}
H^0 \bigl (\sh_{I^{v}(\oo),\bar s_v},j_!^{=h+\delta} HT_\xi(\chi_{v},\Pi_h \{\sfrac{-\delta}{2}\}) \times
\speh_{\delta} (\chi_{v}\{t/2\}))_{\overline \Zm_l} \otimes \Xi^{\delta/2} \bigr)_{\mathfrak m}\\
\to H^0 \bigl (\sh_{I^{v}(\oo),\bar s_v},
j_!^{=h+\delta-1} HT_\xi(\chi_{v},\Pi_h \{\psfrac{1-\delta}{2}\}) \\ \times
\speh_{\delta-1}(\chi_{v}\{t/2\}))_{\overline \Zm_l} \otimes \Xi^{\psfrac{\delta-1}{2}} \bigr)_{\mathfrak m}
\end{multline}
are strict. Then from the previous fact stressed after \eqref{eq-map3}, this property
remains true when we consider the associated spectral sequence for $1 \leq h' \leq h_0(I^v,\chi_v)$.

Consider now $h=h_0(I^v,\chi_v)$,
where we know the torsion to be nontrivial. From what was observed above
we then deduce that the map
\begin{multline} \label{eq-map1-coho}
H^0 \Bigl (\sh_{I^{v}(\oo),\bar s_v},\\
\shoveright{j_!^{=h_0(I^v,\chi_v)+1} HT_\xi\bigl(\chi_{v},\Pi_{h_0(I^v,\chi_v)} \{\sfrac{-1}{2}\}\bigr)
\times \chi_{v}\bigl\{h_0(I^v,\chi_v)/2\bigl\}\bigr)_{\overline \Zm_l} \otimes \Xi^{\sfrac{1}{2}} \Bigr)_{\mathfrak m}} \\
\to H^0 \bigl (\sh_{I^{v}(\oo),\bar s_v},
j_!^{=h_0(I^v,\chi_v)} HT_\xi(\chi_{v},\Pi_{h_0(I^v,\chi_v)})_{\overline \Zm_l} \bigr)_{\mathfrak m}
\end{multline}
has a nontrivial torsion cokernel so that $i_I(h_0(I^v,\chi_v))=0$.

Finally for any $1 \leq h \leq h_0(I^v,\chi_v)$, the map like \eqref{eq-map1-coho} for $h+\delta-1 < h_0(I^v,\chi_v)$
are strict so that the $H^i(\sh_{I^{v}(\oo),\bar s_v},\lexp p j^{=h}_{!*}
HT_\xi(\chi_{v},\Pi_h)_{\overline \Zm_l})_{\mathfrak m}$ are zero for $i < h-h_0$ while when
$h+\delta-1=h_0$ its cokernel has nontrivial torsion, which gives
then a nontrivial torsion class in $H^{h-h_0}(\sh_{I^{v}(\oo),\bar s_v},\lexp p j^{=h}_{!*}
HT_\xi(\chi_{v},\Pi_h)_{\overline \Zm_l})_{\mathfrak m}$.
\end{proof}

\begin{lemma} \label{lem-h0I}
The integers $h_0(I^v,\chi_v)$ and $i_{I^v}(h,\chi_v)$ only depend on the reduction modulo $l$ of $\chi_v$.
\end{lemma}

\begin{proof}
For $P$ a torsion-free $\overline \Zm_l$-perverse sheaf,
recall the well-known short exact sequence
\[
0 \to H^i(X,P) \otimes_{\overline \Zm_l} \overline \Fm_l \to
H^i(X,P \otimes^\Lm_{\overline \Zm_l} \overline \Fm_l) \to
H^{i+1}(X,P)[l] \to 0.
\]
We apply it to $X=\sh_{I^v,\bar s_v}$ and $P=P_\xi(h_0(I^v,\chi_v),\chi_{v})_{\overline \Zm_l}$.
Recall that, thanks to our hypothesis on $\mathfrak m$,
the $\overline \Qm_l$-cohomology of $P$, localized at $\mathfrak m$,
is concentrated in degree $0$.
The same is true for $P_\xi(h_0(I^v,\chi_v),\chi'_{v})_{\overline \Zm_l}$ for any character
$\chi'_v \equiv \chi_v \bmod l$. Recall also that $P_\xi(h_0(I^v,\chi_v),\chi_{v})_{\overline \Zm_l}$
(\resp $P_\xi(h_0(I^v,\chi_v),\chi'_{v})_{\overline \Zm_l}$) is a local system on the stratum
$\sh_{I^v,\bar s_v}^{\geq h_0(I^v,\chi_v)}$ so that
\[
P_\xi(h_0(I^v,\chi_v),\chi_{v})_{\overline \Zm_l} \otimes^\Lm_{\overline \Zm_l} \overline \Fm_l
\simeq P_\xi(h_0(I^v,\chi_v),\chi'_{v})_{\overline \Zm_l} \otimes^\Lm_{\overline \Zm_l} \overline \Fm_l
\]
is simply the reduction modulo $l$ of this local system.
From the definition of $h_0(I^v,\chi_v)$, the cohomology of
$P_\xi(h_0(I^v,\chi_v),\chi_{v})_{\overline \Zm_l}$ has torsion in degrees $0$ and $1$
so that its $\overline \Fm_l$-cohomology, localized at $\mathfrak m$, is
concentrated in degrees $-1,0,1$: moreover, for $h> h_0(I^v,\chi_v)$ the
$\overline \Fm_l$-cohomology of $P_\xi(h,\chi_v)_{\overline \Zm_l}$ is concentrated in degree $0$.

The same is then true for the $\overline \Fm_l$-cohomology,
localized at $\mathfrak m$, of $P_\xi(,h_0(I^v,\chi_v),\chi'_{v})_{\overline \Zm_l}$
so that its $\overline \Zm_l$-cohomology, localized at $\mathfrak m$, must have
torsion in degrees $0$ and $1$. Moreover, for $h>h_0(I^v,\chi_v)$,
the $\overline \Fm_l$-cohomology of $P_\xi(h,\chi'_v)_{\overline \Zm_l}$ is concentrated
in degree $0$. By definition we then have $h_0(I^v,\chi'_v)=h_0(I^v,\chi'_v)$.

Concerning $i_{I^v}(h,\chi'_v)$ the result then follows from the previous lemma.
\end{proof}

From the main assumption of Section \ref{para-max} and as explained in the introduction of
this section, we focus on $\Psi_{\xi,\varrho}$ when $\varrho=\mathds 1_v$
is the trivial character. We are then interested in the characters $\chi_v$
congruent to the trivial character $\mathds 1_v$ modulo $l$.

\begin{nota} \label{nota-h0I}
Following the notation of Proposition \ref{prop-h0I}, we will
denote $h_0(I^v)$ (\resp $i_{I^v}(h)$) for $h_0(I^v,\mathds 1_v)$
(\resp $i_{I^v}(h,\mathds 1_v)$).
\end{nota}

\begin{lemma} \label{lem-iIv}
With the notation of Proposition \ref{prop-h0I}, we have $h_0(I^v) \geq h_0(I)$.
\end{lemma}

\begin{proof}
Consider the previous map \eqref{eq-map1-coho} by replacing $h_0(I^v)$ by $h_0(I)$.
As by hypothesis the order of $q_v$ modulo $l$ is strictly greater than $d$,
then the
pro-order of the local component $I_v$ of $I$ at $v$
is invertible modulo $l$, so that the functor of invariants under $I_v$ is exact.
Note then that, as the $I_v$-invariants of the map \eqref{eq-map1-coho} when
replacing $h_0(I^v)$ by $h_0(I)$, has a cokernel which is not free, then the cokernel
of \eqref{eq-map1-coho}, for $h_0(I)$, is also not free.
\end{proof}

From the previous proof, we also deduce that all cohomology classes of any of
the $H^i(\sh_{I^{v}(\oo),\bar s_v},P_\xi(t,\chi_v)_{\overline \Zm_l})_{\mathfrak m}$ come
from the non strictness of some map \eqref{eq-map1-coho} with $\Pi_v:=\st_t(\chi_v)$. In the following we will focus
on $H^i(\sh_{I^{v}(\oo),\bar s_v},P_\xi(t,\chi_v)_{\overline \Zm_l})_{\mathfrak m}[l]$ as an $\overline \Fm_l$-representation of $\GL_d(F_v)$.
More precisely we are interested in irreducible such subquotients which have
maximal non-degeneracy level at $v$.

\begin{nota}
Fix such a non degeneracy level $\underline{\lambda}$ for $\GL_d(F_v)$
in the sense of Notation \ref{nota-nondegeneracy}, which is maximal for torsion classes
in $H^0(\sh_{I^{v}(\oo),\bar s_v},P_\xi(t,\chi_v)_{\overline \Zm_l})_{\mathfrak m}[l]$ for various $t$ such that $1 \leq t \leq d$ and $\chi_v \in \cusp_{\mathds 1_v}(-1)$.
\end{nota}

\begin{lemma} \label{lem-rem-after1}
Let $\chi_v \in \cusp_{-1}(\mathds 1_v)$.
Then all $\overline \Fm_l[\GL_d(F_v)]$-irreducible subquotients of
$H^i(\sh_{I^{v}(\oo),\bar s_v},P_\xi(t,\chi_v)_{\overline \Zm_l})_{\mathfrak m}[l]$, for $i\neq 0,1$,
have a level of
non degeneracy strictly less than $\underline{\lambda}$.
\end{lemma}

\begin{proof}
It easily follows from the observation that the level of non degeneracy of
the reduction modulo $l$ of $\speh_h(\chi_v) \simeq \chi_v$ is strictly less than
those of the reduction modulo $l$ of $\st_h(\chi_v)$ which, \cf \cite{boyer-repmodl},
is irreducible as the order
of $q_v$ modulo $l$ is strictly greater than $d$ and so strictly greater than $h$.
\end{proof}

\subsection{Global torsion and genericity}
\label{para-generic}

Recall that $v \in \spl$ is such that the order of~$q_v$ modulo $l$ is strictly
greater than $d$. Let us denote by $I^{v}$ the component of $I$ outside $v$.
We then simply denote by $\Psi_v$ and $\Psi_{v,\xi}$,
the inductive system of perverse sheaves indexed by the finite level $I^{v} I_v \in \IC$
for varying $I_v$.

For $\pi_{v} \in \cusp_\varrho$, let us denote by $\Fil^1_{!,\pi_{v}}(\Psi_\varrho)$ the quotient of $\Fil^1_!(\Psi_\varrho)$ such that
$\Fil^1_{!,\pi_{v}}(\Psi_v) \otimes_{\overline \Zm_l} \overline \Qm_l \simeq \Fil^1_!(\Psi_{\pi_{v}})$
where $\Psi_{\pi_{v}}$ is the direct factor of $\Psi_v \otimes_{\overline \Zm_l} \overline \Qm_l$
associated to $\pi_{v}$, \cf \cite{boyer-torsion}.

\begin{remark}
In the following, we will be mainly concerned with the case where $\pi_v$
is a character $\chi_v$ whose reduction modulo $l$ is the trivial character.
We will then write the main statement in this case.
\end{remark}

Recall the following resolution of $\Fil^1_{!,\chi_{v}}(\Psi_\varrho)$
\begin{multline} \label{eq-resolution-psi0}
0 \to j^{=d}_! HT(\chi_{v},\speh_d(\chi_{v}))_{\overline \Zm_l} \otimes \Lm (\chi_{v}(\psfrac{d-1}{2}))\\
\to j^{=d-1}_!HT(\chi_{v},\speh_{d-1}(\chi_{v}))_{\overline \Zm_l} \otimes \Lm (\chi_{v}(\psfrac{d-2}{2}))
\to \cdots \\ \to j^{=1}_! HT(\chi_{v},\chi_{v})_{\overline \Zm_l} \otimes \Lm(\chi_{v})
\to \Fil^1_{!,\chi_{v}}(\Psi_\varrho) \to 0,
\end{multline}
which is proved in \cite{boyer-torsion} over $\overline \Qm_l$. I claim it is
also true over $\overline \Zm_l$. Indeed, using Lemma \ref{lem-ext0},
it is equivalent to the fact the sheaf cohomology of $\Fil^1_{!,\chi_{v}}(\Psi_\varrho)$
is torsion-free, which follows then from \cite{s-s}, the comparison
theorem of Faltings-Fargues \cf \cite{fargues-faltings} and the main theorem
of \cite{fargues-monodromie}.

\begin{remark}
In \cite{boyer-duke}, we prove the same resolution for any irreducible
cuspidal representation $\pi_v$ in place of $\chi_v$.
\end{remark}

More generally for $\gr^t_!(\Psi_\varrho) \onto
\gr^t_{!,\chi_v}(\Psi_\varrho)$, we have
\begin{multline} \label{eq-resolution-psi}
0 \to j^{=d}_! HT(\chi_{v},LT_{\chi_v}(t-1,d-t))_{\overline \Zm_l} \otimes \Lm(\chi_{v}(\psfrac{d-2t+1}{2}))\\
\to j^{=d-1}_!HT(\chi_{v},LT_{\chi_{v}}(t-1,d-t-1))_{\overline \Zm_l} \otimes \Lm(\chi_{v}(\psfrac{d-2t}{2}))
\to \cdots \\ \to j^{=t}_! HT(\chi_{v},\st_t(\chi_{v}))_{\overline \Zm_l} \otimes \Lm(\chi_{v})
\to \gr^t_{!,\chi_{v}}(\Psi_\varrho) \to 0.
\end{multline}
Finally all the torsion cohomology classes of the
$H^i(\sh_{I^v(\oo),\bar s_v},\gr^t_{!,\chi_v}(\Psi_\varrho))_{\mathfrak m}$
come from the non strictness of the maps
\begin{multline} \label{eq-map-hmax}
H^{0}(\sh_{I^v(\oo),\bar s_v},j^{=h+1}_! HT(\chi_v,\Pi_{h+1})_{\overline \Zm_l})_{\mathfrak m}\\
\to H^0(\sh_{I^v(\oo),\bar s_v},j^{=h}_! HT(\chi_v,\Pi_h)_{\overline \Zm_l})_{\mathfrak m},
\end{multline}
where $(\Pi_h,\Pi_{h+1})$ is of the shape $\bigl (LT_{\chi_v}(t-1,h-t),LT_{\chi_v}(t-1,h+1-t) \bigr)$.

We can then copy the proof of Lemma \ref{lem-torsion-min} which gives us the following
statement.

\begin{lemma}
For every $h$ such that $1 \leq h \leq h_0(I^v)$, the number
$i_I(h)=h-\nobreak h_0(I^v)$ of Notation \ref{nota-ih}, is also the lowest integer $i$ such that the torsion
of $H^i(\sh_{I^{v}(\oo),\bar s_v},\gr^h_{!,\chi_{v}}(\Psi_{\varrho,\xi}))_{\mathfrak m}$ is non zero.
\end{lemma}

\begin{lemma} \label{lem-generic-sq}
For every $i$, the $l$-torsion of
\[
H^i(\sh_{I^v(\oo),\bar \eta_v},V_{\xi,\overline \Zm_l})_{\mathfrak m},
\]
as an $\overline \Fm_l[\GL_d(F_v)]$-module,
does not have an irreducible generic subquotient
whose cuspidal support is made of characters.
\end{lemma}

\begin{remark}
Note that when the order of $q_v$ modulo $l$ is strictly greater than $d$, then
there is no difference between cuspidal or supercuspidal support made of characters.
\end{remark}

\begin{proof}
Recall first that, as by hypothesis $\overline{\rho_{\mathfrak m}}$ is irreducible, the
$\overline \Qm_l$-version of the spectral sequence \eqref{eq-ss2} degenerates at $E_1$ so
that in particular all the torsion cohomology classes appear in the $E_1$ terms.
As we are only interested in representations with cuspidal support made of characters,
we only have to deal with the perverse sheaves $P(t,\chi_v)_{\overline \Zm_l}$ so that the result
follows from the previous maps \eqref{eq-map-hmax} and the fact that for
any $r>0$, the reduction modulo $l$ of $LT_{\chi_v}(t-1,r)$ does not admit any irreducible
generic subquotient.
\end{proof}

\subsection{Torsion and modified lattices}
\label{para-modification}

Recall that we argue by contradiction, assuming there exists a finite level $I$
unramified at the place $v$, such that the torsion of some of
the $H^i(\sh_{I,\bar \eta_v},V_{\xi,\overline \Zm_l})_{\mathfrak m}$ is non
zero. We then increase the level at $v$ to infinity and define the index
$h_0(I^v)$ which might be greater than the index $h_0(I)$ defined in level $I$,
\cf Notation \ref{nota-h0I}, of the first Harris-Taylor perverse sheaf
associated to a character $\chi_v$ congruent to $\mathds 1_v$ modulo $l$,
with nontrivial torsion cohomology class.

We now come back to a finite level at $v$ with two main objectives:
first we want to keep the torsion in the cohomology of
$P_\xi(\chi_v,h_0)_{\overline \Zm_l}$ and secondly we intend to simplify the spectral sequence
of vanishing cycles.

\begin{remark}
The main reason to go to infinite level at $v$ is to be able to use
the notion of level of non degeneracy.

To be able to deal with representations, we fix a place $w \neq v$
with $w \in \spl(I)$ and verifying the
same hypothesis as $v$, \ie $q_w$ modulo $l$ is of order strictly greater than~$d$.
\end{remark}

\begin{nota}
We then denote as before by $I^w(\oo)$ when the level is infinite at~$w$ and
$I^{w,v}(\oo)$ when the level is infinite at $v$ and $w$. We also denote $h_0$
for $h_0(I^w)$, the highest index when torsion appear in the cohomology of
a Harris-Taylor perverse sheaf in infinite level at $w$ and maximal level at $v$.
\end{nota}

\begin{lemma} \label{lem-oo-finite}
With the notations of \ref{lem-iwahori},
$H^0(\sh_{I^{w,v}(\oo),\bar s_v},P_\xi(h_0,\mathds 1_v)_{\overline \Zm_l})_{\mathfrak m}$ has nontrivial torsion
classes invariant under $\Iw_v(d-h_0,1,\dots,1)$.
\end{lemma}

\begin{proof}
Note first that by our absurd hypothesis in Section \ref{para-max},
\[
H^0(\sh_{I^{w,v}(\oo),\bar s_v},j^{=h_0}_{!*}
HT_\xi(\mathds 1_v,\speh_{h_0}(\mathds 1_v))_{\overline \Zm_l})_{\mathfrak m}
\]
has
nontrivial torsion classes invariant under the action of $\GL_d(\OC_v)$ coming
from the map \eqref{eq-map1-coho} by putting
$\speh_{h_0+1}(\mathds 1_v)$ (\resp $\speh_{h_0}(\mathds 1_v)$)
in place of
\[
\Pi_{h_0(I^v,\chi_v)} \{\sfrac{-1}{2}\} \times \chi_v \{h_0(I^v,\chi_v)/2\}
\]
(\resp $\Pi_{h_0(I^v,\chi_v)}$).
Then by putting $\st_{h_0}(\mathds 1_v)$ in place of $\Pi_{h_0(I^v,\chi_v)}$
in \eqref{eq-map1-coho}, we deduce that
$H^0(\sh_{I^{w,v}(\oo),\bar s_v},P_\xi(h_0,\mathds 1_v))_{\mathfrak m}$ has nontrivial torsion
classes invariant under $\Iw_v(d-h_0,1,\dots,1)$, where, with
the notations of \ref{lem-iwahori}, $(d-h_0,1,\dots,1)$ is the dual
partition of $(h_0+1,1,\dots,1)$.
\end{proof}

\begin{nota}
We will now consider the following level
\[
I^w(h_0):=I^w(\oo) \Iw_v(d-h_0,1,\dots,1),
\]
which is infinite at $w$ and of Iwahori type at $v$.
\end{nota}

\begin{remark}
Dealing with infinite level at $w$ allows to talk about
representations of $\GL_d(F_w)$, while
Iwahori type subgroup allows to simplify the spectral sequence. Indeed
from Lemma \ref{lem-iwahori}, and using the definition of $LT_{\chi_v}(t,s)$
through an induced representation, \cf Definition \ref{defi-rep},
we note that
for $h \geq h_0+2$, then $LT_{\chi_v}(h,d-h-1)$ does not have any nontrivial
vector invariant by $\Iw_v(d-h_0,1,\dots,1)$. Moreover for an irreducible
representation $\pi_v$ of $\GL_{d-h}(F_v)$ with $h \geq h_0+1$,
then \hbox{$LT_{\chi_v}(h_0,h-h_0) \times \pi_v$}
admits nontrivial invariant vectors by $\Iw_v(d-h_0,1,\dots,1)$ if and only if
$\pi_v \simeq \chi_{v,1} \times \cdots \times \chi_{v,d-h}$ is unramified, \cf the
examples following Lemma \ref{lem-iwahori}.
Similar simplifications also appear in Lemma \ref{lem-facts0}.
\end{remark}

We then focus on the free quotient of
$H^0(\sh_{I^w(h_0),\bar \eta_v},V_{\xi,\overline \Zm_l}[d-1])_{\mathfrak m}$
by means of the spectral sequence of vanishing cycles. From \eqref{eq-psi-dec},
we are then lead to study the cohomology of $\Psi_{\xi,\mathds 1_v}$
by means of its filtration of stratification and so we first focus
on the cohomology of $\gr^{h_0}_{!,\chi_v}(\Psi_{\xi,\mathds 1_v})$ for a character
$\chi_v \in \cusp_{-1}(\mathds 1_v)$ which is by definition a quotient of
$\gr^{h_0}_{!}(\Psi_{\xi,\mathds 1_v})$.
To do so, consider first the filtration constructed in \cite{boyer-torsion}:
\[
\Fil^{d-h_0} (\gr^{h_0}_{!,\chi_v}(\Psi_{\xi,\mathds 1_v})) \subset
\cdots \subset \Fil^0(\gr^{h_0}_{!,\chi_v}(\Psi_{\xi,\mathds 1_v})),
\]
with successive free graded parts
$\gr^i(\gr^{h_0}_{!,\chi_v}(\Psi_{\xi,\mathds 1_v}))$ which, for $i \geq 0$,
is a $\overline \Zm_l$-structure of
the $\overline \Qm_l$-perverse sheaf $P(h_0+i,\chi_v)_{\overline \Zm_l}(\psfrac{1-h_0+i}{2})$.
We then introduce two $\overline \Zm_l$-lattices of
\begin{equation} \label{eq-def-sq}
H^0(\sh_{I^w(h_0),\bar s_v},P_\xi(h_0+i,\chi_v)_{\overline \Zm_l})_{\mathfrak m}
\otimes_{\overline \Zm_l} \overline \Qm_l.
\end{equation}
\begin{itemize}
\item The first one denoted by $\Gamma_{\xi,\chi_v,\mathfrak m}(I,h_0+i)$ is given by the free $\overline \Zm_l$-cohomology: recall that as the order of $q_v$
modulo $l$
is supposed to be strictly greater than $d$ then, \cf \cite{boyer-repmodl},
the reduction modulo $l$ of
$\st_{h_0+i}(\chi_v)$ and that of $\chi_v[t]_D$,
remains irreducible so that, up to homothety, there is a unique
stable lattice of $P_\xi(h_0+i,\chi_v)_{\overline \Zm_l}$.

\item The spectral sequence associated to the previous filtration of
$\gr^{h_0}_{!,\chi_v}(\Psi_{\xi,\mathds 1_v})$ provides a filtration of $H^0_{\free}
(\sh_{I^w(h_0),\bar s_v},\gr^{h_0}_{!,\chi_v}(\Psi_{\xi,\mathds 1_v}))_{\mathfrak m}$,
and $\Gamma_{\xi,\chi_v,!,\mathfrak m}(I,h_0+i,h_0)$ is then the lattice of the
subquotient in this filtration corresponding to \eqref{eq-def-sq}.
\end{itemize}

By construction we have\vspace*{-3pt}
\begin{equation} \label{eq-T}
\Gamma_{\xi,\chi_v,\mathfrak m}(I,h_0+1)
\hto \Gamma_{\xi,\chi_v,!,\mathfrak m}(I,h_0+1,h_0),
\end{equation}
but the cokernel of torsion might be nontrivial due to torsion
in the remaining of the $E_\oo$ terms. The main point is first to show,
under the absurd hypothesis of Section~\ref{para-max}, that this cokernel is nontrivial and
then, to prove some property verified by its $l$\nobreakdash-tor\-sion and finally achieve to a
contradiction. We then first focus on the first step about nontriviality.

\begin{lemma} \label{lem-important}
The cokernel $T$ of \eqref{eq-T}:\vspace*{-3pt}
\[
0 \to \Gamma_{\xi,\chi_v,\mathfrak m}(I,h_0+1)
\to \Gamma_{\xi,\chi_v,!,\mathfrak m}
(I,h_0+1,h_0) \to T \to 0,
\]
is non zero and every irreducible subquotient of its $l$-torsion as a
$(\Tm_{\xi,\mathfrak m}^S \times\nobreak \GL_d(F_w)) \otimes_{\overline \Zm_l} \overline \Fm_l$-module,
can be obtained as a subquotient of the torsion submodule of
the cokernel~of\vspace*{-3pt}
\begin{multline} \label{eq-torsion-trop2}
H^{0}(\sh_{I^w(h_0),\bar s_v},j^{=h_0+1}_!
HT_\xi(\chi_v,\st_{h_0+1}(\chi_v))_{\overline \Zm_l})_{\mathfrak m} \\
\to H^0(\sh_{I^w(h_0),\bar s_v},j^{=h_0}_!
HT_\xi(\chi_v,\st_{h_0}(\chi_v))_{\overline \Zm_l})_{\mathfrak m}.
\end{multline}
\end{lemma}

\begin{proof}
The idea is to compute the cohomology of $\gr^{h_0}_{!}(\Psi_{\xi,\chi_v})$
in two different ways, first by means of the spectral sequence associated to
\eqref{eq-resolution-psi} and secondly by means of its filtration of stratification with graded parts the Harris-Taylor perverse sheaves.

To argue we will rest on the level of non degeneracy at $v$ so that we pass to
$I^{w,v}(\oo)$-level: as $q_v$ modulo $l$ is of order $>d$ taking
invariant under any sub-group of $\GL_d(\OC_v)$ is an exact functor and it
will be easy to go down to level $I^w(h_0)$.

First note that the $I^{w,v}(\oo)$-version of
\eqref{eq-torsion-trop2} is non strict
if and only if the same is true for its non induced version in the next formula,
whatever a representation $\Pi_{h_0}$ of $\GL_{h_0}(F_v)$ is:\vspace*{-3pt}
\begin{multline} \label{eq-torsion-trop3}
H^{0}(\sh_{I^{w,v}(\oo),\bar s_v,\overline{1_{h_0}}},j^{=h_0+1}_{\overline{1_{h_0}},!} HT_{\xi}
(\chi_v,\Pi_{h_0} \otimes \chi_v)_{\overline{1_{h_0}},\overline \Zm_l})_{\mathfrak m} \\
\to H^0(\sh_{I^{w,v}(\oo),\bar s_v,\overline{1_{h_0}}},j^{=h_0}_{\overline{1_{h_0}},!}
HT_{\xi}(\chi_v,\Pi_{h_0})_{\overline{1_{h_0}},\overline \Zm_l})_{\mathfrak m},
\end{multline}
where we set
\[
j^{=h_0+1}_{\overline{1_{h_0}}}:
\sh^{=h_0+1}_{I^v(\oo),\bar s_v,\overline{1_{h_0}}}
\hto \sh^{\geq 1}_{I^v(\oo),\bar s_v},
\]
where $\sh^{=h_0+1}_{I^{w,v}(\oo),\bar s_v,\overline{1_{h_0}}}$ is the
disjoint union of the pure strata, \cf Notation \ref{nota-strata},
$\sh^{=h_0+1}_{I^{w,v}(\oo),\bar s_v,g}$
contained in $\sh^{\geq h_0}_{I^{w,v}(\oo),\bar s_v,\overline{1_{h_0}}}$.

\skpt
\begin{lemma} \label{lem-facts0}
\begin{enumeratei}
\item\label{lem-facts0i} With $\Pi_{h_0}=\st_{h_0}(\chi_v)$, the cokernel of the induced version of
\eqref{eq-torsion-trop3} has nontrivial vectors invariant
under $\Iw_v(d-h_0,1,\dots,1)$.

\item\label{lem-facts0ii} For $h > h_0$, whatever are the representations $\Pi_h$ and
$\Pi_{h+1}$ of respectively $\GL_h(F_v)$ and $\GL_{h+1}(F_v)$, the cokernel of
\begin{multline} \label{eq-torsion-trop-h}
H^{0}(\sh_{I^{w,v}(\oo),\bar s_v},j^{=h+1}_!
HT_\xi(\chi_v,\Pi_{h+1})_{\overline \Zm_l})_{\mathfrak m} \\
\to H^0(\sh_{I^{w,v}(\oo),\bar s_v},j^{=h}_!
HT_\xi(\chi_v,\Pi_h)_{\overline \Zm_l})_{\mathfrak m},
\end{multline}
does not have any non zero invariant vector under $\GL_d(\OC_v)$.

\item\label{lem-facts0iii} For $\Pi_h=LT_{\chi_v}(h_0-1,h-h_0)$ and
$\Pi_{h+1}=LT_{\chi_v}(h_0-1,h-h_0+1)$, as in \eqref{eq-map-hmax},
the cokernel of \eqref{eq-torsion-trop-h} does not have any non zero invariant vector under
$\Iw_v(d-h_0,1,\dots,1)$.
\end{enumeratei}
\end{lemma}

\begin{proof}
The integer $h_0$ is chosen so that, by imposing $\Pi_{h_0}$ to be unramified,
using also the fact that $q_v$ modulo $l$ is of order $>d$ so that the functor of $P_{h_0,d}(\OC_v)$-invariants is exact,
then the cokernel of \eqref{eq-torsion-trop3}
has nontrivial vectors invariant under $P_{h_0,d}(\OC_v)$.
Then when $\Pi_{h_0}=\st_{h_0}(\chi_v)$, by modifying the factor
$\GL_{d-h_0}(\OC_v)$ by its classical Iwahori subgroup, we then deduce \eqref{lem-facts0i}.

\eqref{lem-facts0ii} It follows from the definition of $h_0$ and the fact that
the functor of $\GL_d(\OC_v)$-invariants is exact.

\eqref{lem-facts0iii} If there were nontrivial zero invariant vectors under $\Iw_v(d-h_0,1,\dots,1)$
then replacing $\Pi_h$ and $\Pi_{h+1}$ respectively by $\speh_h(\chi_v)$ and
$\speh_{h+1}(\chi_v)$, there would exist non zero invariants under
$\GL_d(\OC_v)$, which is not the case by \eqref{lem-facts0ii}, \cf the examples following
Lemma \ref{lem-iwahori}.
\end{proof}

We then compute the $\mathfrak m$-localized cohomology of
$\gr^{h_0}_{!,\chi_v}(\Psi_{\xi,\mathds 1_v})$
in level $I^{w,v}(\oo)$ having nontrivial invariant under
$\Iw_v(d-h_0,1,\dots,1)$.
By maximality of $h_0$, note that for $h_0<t \leq d$, the cohomology groups of
\[
P_\xi(t,\chi_v)_{\overline \Zm_l}\quad\text{and}\quad j^{=t}_! HT_\xi(\pi_v,LT_{\chi_v}(h_0-1,t-h_0))_{\overline \Zm_l},
\]
after localization by $\mathfrak m$, do no
have any nontrivial torsion vector invariant under \hbox{$\Iw_v(d-h_0,1,\dots,1)$}
as explained in the previous lemma.

(1) Following the proof of \ref{lem-torsion-min} with the spectral sequence
associated to
\eqref{eq-resolution-psi} and neglecting torsion classes
which do not have nontrivial vectors invariant by \hbox{$\Iw_v(d-h_0,1,\dots,1)$},
we then deduce that
$H^i_{\tor}(\sh_{I^{w,v}(\oo),\bar s_v},
\gr^{h_0}_{!}(\Psi_{\xi,\chi_v}))_{\mathfrak m}$ does not have any nontrivial vector invariant under $\Iw_v(d-h_0,1,\dots,1)$ if $i \neq 0,1$,
while for $i=0$ the torsion is nontrivial and the
vectors invariant by $\Iw_v(d-h_0,1,\dots,1)$ are given by the non strictness of
\begin{multline} \label{eq-torsion-trop}
H^{0}(\sh_{I^{w,v}(\oo),\bar s_v},j^{=h_0+1}_! HT_\xi(\chi_v,LT_{\chi_v}(h_0-1,1))_{\overline \Zm_l})_{\mathfrak m} \\
\to H^0(\sh_{I^{w,v}(\oo),\bar s_v},j^{=h_0}_!
HT_\xi(\chi_v,\st_{h_0}(\chi_v))_{\overline \Zm_l})_{\mathfrak m}.
\end{multline}

(2) Concerning $H^0(\sh_{I^{w,v}(\oo),\bar s_v},
P_\xi(h_0,\chi_v)_{\overline \Zm_l})_{\mathfrak m}$,
its torsion submodule is
parabolically induced, so that beside those coming from the non strictness of \eqref{eq-torsion-trop},
there is also the contribution given by the non strictness of \eqref{eq-torsion-trop2}, which
contains in particular a subquotient, denoted $\widetilde T$,
such that $\widetilde T[l]$ is of level of non degeneracy strictly greater than those appearing
in \eqref{eq-torsion-trop}. Note moreover that
\begin{itemize}
\item $\widetilde T[l]$ has nontrivial vectors under $\Iw_v(d-h_0,1,\dots,1)$;

\item as $(\Tm^S_{\xi,\mathfrak m} \times \GL_d(F_w)) \otimes_{\overline \Zm_l} \overline \Fm_l$-modules, the irreducible subquotients of
$\widetilde T[l]$ are also subquotients of the torsion of the $l$-torsion of the
cokernel of \eqref{eq-torsion-trop}. Indeed, $\widetilde T[l]$ is given as the cokernel
of \eqref{eq-torsion-trop} where we replace $LT_{\chi_v}(h_0-1,1)$ by
$\st_{h_0+1}(\chi_v)$.
\end{itemize}

(3) Consider then the cohomology of $\gr^{h_0}_{!,\chi_v}(\Psi_{\varrho,\xi})$
computed through
its filtration of stratification with graded parts, up to Galois shifts,
\[
H^0(\sh_{I^{w,v}(\oo),\bar s_v},P_\xi(h_0+k,\chi_v)_{\overline \Zm_l})_{\mathfrak m},
\]
for $0 \leq k \leq d-h_0$,
and more particularly \emph{the induced filtration of the free quotient of}
$H^0(\sh_{I^{w,v}(\oo),\bar s_v},\gr^{h_0}_{!,\chi_v}(\Psi_{\mathds 1_v,\xi}))_{\mathfrak m}$ as before.
As the level of non degeneracy of $\widetilde T[l]$ is higher than that of the $l$-torsion
of $H^0(\sh_{I^{w,v}(\oo),\bar s_v},\Fil^{h_0}_{!,\chi_v}
(\Psi_{\mathds 1_v,\xi}))_{\mathfrak m}$,
computed by means of the spectral sequence associated to \eqref{eq-resolution-psi},
they must be graded parts of this filtration.
We then have a filtration of the free quotient of $H^0(\sh_{I^{w,v}(\oo),\bar s_v},
\gr^{h_0}_{!,\chi_v}(\Psi_{\mathds 1_v,\xi}))_{\mathfrak m}$ for which, among the graded parts, appear
\begin{itemize}
\item torsion modules such as $\widetilde T$,

\item and the free graded parts which are lattices
$\Gamma_{\xi,\chi_v,\mathfrak m}(I^v,h_0+i)$
of the free quotient of the localized cohomology of $P_\xi(\chi_v,h_0+i)$
for $0 \leq i \leq d-h_0$.
\end{itemize}
We now go back to the level $I^w(h_0)=I^w(\oo)\Iw_v(d-h_0,1,\dots,1)$:
as $q_v$ modulo $l$ is of order strictly
greater than $d$, the functor of $\Iw_v(d-h_0,1,\dots,1)$-invariants
is exact. As only the
cohomology of $P_\xi(\chi_v,h_0+i)$ for $i=0,1$ contributes, the result follows from
the fact that $\widetilde T$ has a nontrivial invariant vector under $\Iw_v(d-h_0,1,\dots,1)$.
\end{proof}

Recall that
\[
\gr^{h_0}_!(\Psi_{\xi,\mathds 1_v}) \otimes_{\overline \Zm_l} \overline \Qm_l \simeq
\bigoplus_{\chi_v \in \cusp_{\mathds 1_v}} \gr^{h_0}_{!}(\Psi_{\xi,\chi_v})
\]
so that we can find a filtration of $\gr^{h_0}_!(\Psi_{\xi,\mathds 1_v})$ whose graded parts are
free and isomorphic, after tensoring with $\overline \Qm_l$, to
$\gr^{h_0}_{!}(\Psi_{\xi,\chi_v})$.
Arguing as in the proof of Lemma~\ref{lem-torsion-min}, using \eqref{eq-map-hmax},
we have the following result.

\begin{lemma} \label{lem-important2}
For every $t\geq1$, let $j(t)$ be the minimal integer $j$ such that the torsion of
$H^j(\sh_{I^{w,v}(\oo),\bar s_v},\gr^t_{!}(\Psi_{\xi,\mathds 1_v}))_{\mathfrak m}$ has
nontrivial invariant vectors under $\Iw_v(d-h_0,1,\dots,1)$. Then
\[
j(t)=\begin{cases}
+ \oo & \text{if } t \geq h_0+1, \\
t-h_0 & \text{for } 1 \leq t \leq h_0.
\end{cases}
\]
Moreover as a $(\Tm_{\xi,\mathfrak m} \times \GL_d(F_w)) \otimes_{\overline \Zm_l} \overline \Fm_l$-module,
up to multiplicities, the irreducible subquotients of
$H_{\tor}^{j(t)}(\sh_{I^{w}(h_0),\bar s_v},\gr^t_{!}
(\Psi_{\xi,\mathds 1_v}))_{\mathfrak m}$ are
independent of $t$.
\end{lemma}

\begin{prop} \label{prop-important}
Up to multiplicities, the set of irreducible
$\overline \Fm_l[\GL_d(F_w)]$-subquotients of the $l$-torsion of\footnote{or
those of $H^{0}(\sh_{I^w(h_0),\bar s_v},P_\xi(h_0,\chi_v))_{\mathfrak m}$
as explained above}
$H^{0}(\sh_{I^w(h_0),\bar s_v},\gr^{h_0}_{!}(\Psi_{\xi,\mathds 1_v}))_{\mathfrak m}$,
are the same as those of
$H^{d-h_0}(\sh_{I^w(h_0),\bar s_v},V_{\xi,\overline \Zm_l})_{\mathfrak m}$.
\end{prop}

\begin{proof}
We compute $H^{d-h_0(I^v)}(\sh_{I^w(h_0),\bar s_v},V_{\xi,\overline \Zm_l})_{\mathfrak m}$
using the filtration $\Fil^{\csbullet}_!(\Psi_{\xi,\mathds 1_v})$ through the spectral sequence
\eqref{eq-ss2}. Recall that for every $p+q \neq 0$, the free quotient of $E^{p,q}_{!,\varrho,1}$
are zero. By definition of the filtration, these $E^{p,q}_{!,\varrho,1}$ are trivial for $p \geq 0$
while, thanks to the previous lemma, for any $p \leq -1$ they are zero for
$p+q < j(p):=p-h_0$. Note then that
$E_{!,\varrho,1}^{-1,j(1)+1}$, which is torsion and non zero, according to the previous lemma,
is equal to $E_{!,\varrho,\oo}^{j(1)} \simeq
H^{d-h_0}(\sh_{I^w(h_0),\bar s_v},V_{\xi,\overline \Zm_l})_{\mathfrak m}$.
\end{proof}

Combining the result of Lemma \ref{lem-important} with the previous
proposition, we then deduce that the cokernel $T$ of Lemma \ref{lem-important} verifies
the following property. As an $\overline \Fm_l$\nobreakdash-representation of $\GL_d(F_w)$, every irreducible
subquotient of $T[l]$ is also a subquotient of
$H^{d-h_0}(\sh_{I^w(h_0),\bar s_v},V_{\xi,\overline \Zm_l})_{\mathfrak m}$.
Then applying Lemma \ref{lem-generic-sq} at the place $w$, which satisfies the same hypothesis as $v$, we then deduce the following result.

\begin{corol} \label{coro-absurd}
As an $\overline \Fm_l$-representation of $\GL_d(F_w)$, the $l$-torsion of
the cokernel $T$ of Lemma \ref{lem-important}
does not contain any irreducible generic subquotient with cuspidal support made of characters.
\end{corol}

We can now repeat the arguments with $\gr^k_{!,\chi_v}(\Psi_{\xi,\mathds 1_v})$
for any $1 \leq k \leq h_0$. More precisely, \cf the last remark of Section \ref{para-fil-psi},
consider $\Fil^{i}(\gr^{k}_{!,\chi_c}(\Psi_{\xi,\mathds 1_v}))$
for $i=h_0-k$ and $i=h_0-k+2$. As by hypothesis, the torsion of
$H^{d-1}(\sh_{I,\bar \eta},V_{\xi,\overline \Zm_l})_{\mathfrak m}$
is nontrivial, then there exists $\widetilde{\mathfrak m} \subset \mathfrak m$ such that
\[
\Pi_{\widetilde{\mathfrak m}} \simeq
\st_{h_0+1}(\chi_v) \times \chi_{v,1} \times \cdots \times \chi_{v,d-h_0-1},
\]
with $\chi_v \equiv \mathds 1_v \bmod l$.
Moreover as before
\begin{itemize}
\item $\Fil^{h_0+2-k}(\gr^{k}_{!,\chi_v}(\Psi_{\xi,\mathds 1_v}))$
has trivial cohomology groups in level $I^w(h_0)$ because
the irreducible constituents of
$\Fil^{h_0+2-k}(\gr^{k}_{!,\chi_v}(\Psi_{\xi,\mathds 1_v}))
\otimes_{\overline \Zm_l} \overline \Qm_l$ are, up to Galois shift,
Harris-Taylor perverse
sheaves $P(t,\chi_v)$ with $t \geq h_0+2$;

\item we can apply the previous argument relatively to
$\gr^{h_0}_{!,\chi_v}(\Psi_{\xi,\mathds 1_v})$ to the
quotient
\[
Q:=\Fil^{h_0-k}(\gr^{k}_{!,\chi_v}(\Psi_{\xi,\mathds 1_v}))
/\Fil^{h_0+2-k}(\gr^{h_0}_{!,\chi_v}(\Psi_{\xi,\mathds 1_v})),
\]
so that, denoting by
$\Gamma'_{\xi,\chi_v,!,\mathfrak m} (I,h_0+1,k)$
the lattice of \eqref{eq-def-sq} given by the free quotient of
$H^0(\sh_{I^w(\oo),\bar s_v},Q)_{\mathfrak m}$, the cokernel $T'_k$ of
\[
0 \to \Gamma_{\xi,\chi_v,\mathfrak m}(I,h_0+1)
\to \Gamma'_{\xi,\chi_v,!,\mathfrak m}
(I,h_0+1,k)
\]
is such that $T'_k[l] \neq (0)$ and, as an $\overline \Fm_l$-representation of
$\GL_d(F_w)$, it does not contain any irreducible generic subquotient made of
characters.
\end{itemize}

In addition of the previous arguments, we also have to deal with the
torsion in the cohomology groups of
\[
\gr^{k}_{!,\chi_v}(\Psi_{\xi,\mathds 1_v}) /\Fil^{h_0-k}(\gr^{k}_{!,\chi_v}
(\Psi_{\xi,\mathds 1_v})),
\]
which could modify the lattice $\Gamma'_{\xi,\chi_v,!,\mathfrak m}
(I,h_0+1,k)$ to give the good one denoted above by
$\Gamma_{\xi,\chi_v,!,\mathfrak m}(I,h_0+1,k)$. Note
again that,
as an $\overline \Fm_l$-representation of $\GL_d(F_w)$, this $l$-torsion
does not contain any
irreducible generic subquotient made of characters, so the cokernel of
\[
\Gamma'_{\xi,\chi_v,!,\mathfrak m} (I,h_0+1,k)
\hto \Gamma_{\xi,\chi_v,!,\mathfrak m}(I,h_0+1,k),
\]
is again such that, as an $\overline \Fm_l$-representation of
$\GL_d(F_w)$, its $l$-torsion does not contain any irreducible generic
subquotient made of characters.

Forgetting again Galois shifts, we then conclude that the $l$-torsion of
the cokernel~$T_k$~of
\begin{equation} \label{eq-absurd}
0 \to \Gamma_{\xi,\chi_v,\mathfrak m}(I,h_0+1)
\to \Gamma_{\xi,\chi_v,!,\mathfrak m}
(I,h_0+1,k)
\end{equation}
is non zero and,
as an $\overline \Fm_l$-representation of $\GL_d(F_w)$, it
does not contain any
irreducible generic subquotient made of characters.

\medskip

We now compute $H^{0}(\sh_{I^w(h_0),\bar s_v},\Psi_{\xi,\mathds 1_v}
)_{\mathfrak m}$ by means of the spectral
sequence of vanishing cycles using the filtration
\[
\Fil^{1}_!(\Psi_{\xi,\mathds 1_v}) \harrow \cdots \harrow \Fil^{d-1}_!
(\Psi_{\xi,\mathds 1_v}) \harrow
\Fil^{d}_!(\Psi_{\xi,\mathds 1_v}) \harrow \Psi_{\xi,\mathds 1_v}.
\]
Recall that we can filter each of the
$\gr^k_!(\Psi_{\xi,\mathds 1_v})$
such that the graded parts are, after tensoring with $\overline \Qm_l$ and up to
Galois shift, of the form $P_\xi(\chi_v,t)$ with $\chi_v \in \cusp_{-1}(\mathds 1_v)$.
As before, arguing by contradiction,
we suppose that the torsion of
$H^{d-1}(\sh_{I,\bar \eta_v},V_{\xi,\overline \Zm_l})_{\mathfrak m}$
is nontrivial, and we pay special attention to the lattices of
\begin{equation} \label{eq-def-Vrho}
V_{\xi,\chi_v,\mathfrak m}(I,h_0+1)(\delta):=
H^i(\sh_{I,\bar s_v},P_\xi(h_0+1,\chi_v))_{\overline \Zm_l}(\delta)_{\mathfrak m}
\otimes_{\overline \Zm_l} \overline \Qm_l,
\end{equation}
for $\chi_v \in \cusp_{-1}(\mathds 1_v)$ and various $\delta$.

\begin{enumeratea}
\item
We first start with $\delta=-h_0/2$. Note that by our choice of Iwahori
subgroup,
$H^i(\sh_{I^w(h_0),\bar s_v},\Psi_{\xi,\mathds 1_v}/\Fil^{h_0+1}_!
(\Psi_{\xi,\mathds 1_v}))_{\mathfrak m}$
are all zero. Indeed the torsion-free graded parts $\grr^k(\Psi_{\xi,\mathds 1_v})$
of any exhaustive filtration of
$\Psi_{\xi,\mathds 1_v} / \Fil^{h_0+1}_!(\Psi_{\xi,\mathds 1_v})$, up to Galois torsion, are such that
$\grr^k(\Psi_{\xi,\mathds 1_v}) \otimes_{\overline \Zm_l} \overline \Qm_l \simeq P(t,\chi_v)$
with $\chi_v \in \cusp_{-1}(\mathds 1_v)$ and $t \geq h_0+2$. Then
every irreducible constituent of
$H^i(\sh_{I^{w,v}(\oo),\bar s_v},\grr^k(\Psi_{\xi,\mathds 1_v}))_{\mathfrak m}
\otimes_{\overline \Zm_l} \overline \Fm_l$,
as an $\overline \Fm_l$-representation of $\GL_d(F_v)$ is a subquotient of an induced
representation
$r_l(\st_t(\chi_v))\{\delta/2\} \times \tau$ for some irreducible $\overline \Fm_l$-representation
$\tau$ of $\GL_{d-t}(F_v)$. In~particular such a representation does not have nontrivial invariants under \hbox{$\Iw_v(d-h_0,1,\dots,1)$}, so that, as the functor of
$\Iw_v(d-h_0,1,\dots,1)$-invariants is exact,
there is no cohomology in level $I^w(h_0)$ as stated.

We then deduce that
$H^0(\sh_{I^w(h_0),\bar s_v},P_{\xi,\overline \Qm_l}(h_0+1,\chi_v))(\delta)_{\mathfrak m}$
is a quotient of
$H^0(\sh_{I^w(h_0),\bar s_v},\Psi_{\xi,\mathds 1_v})_{\mathfrak m}
\otimes_{\overline \Zm_l} \overline \Qm_l$
and we denote by
\[
\Gamma_{\xi,\chi_v,\Psi,\mathfrak m}(I,h_0+1,+)
\]
its stable lattice induced by
$H^0_{\free}(\sh_{I^w(h_0),\bar s_v},\Psi_{\xi,\mathds 1_v})_{\mathfrak m}$.

\begin{prop} \label{prop-lattice-psi}
With the previous notations, we have an exact sequence
\[
0 \to \Gamma_{\xi,\chi_v,\Psi,\mathfrak m}(I,h_0+1,+)
\to \Gamma_{\xi,\chi_v,\mathfrak m}(I,h_0+1) \to T
\to 0,
\]
where $T[l]$, as a $\overline{\mathbb F}_l$-representation of $\GL_d(F_w)$,
does not have any irreducible generic subquotient with cuspidal support
made of characters.
\end{prop}

\begin{proof}
We compute $H^0(\sh_{I^w(h_0),\bar s_v},\Psi_{\xi,\mathds 1_v})_{\mathfrak m}$ by means of
the spectral sequence associated to the filtration
\[
\Fil^{1}_!(\Psi_{\xi,\mathds 1_v}) \harrow \cdots \harrow \Fil^{d-1}_!
(\Psi_{\xi,\mathds 1_v})
\harrow \Fil^{d}_!(\Psi_{\xi,\mathds 1_v}) \harrow \Psi_{\xi,\mathds 1_v}.
\]
Recall that $H^0(\sh_{I^w(h_0),\bar s_v},\Psi_{\xi,\mathds 1_v})_{\mathfrak m}=
H^0(\sh_{I^w(h_0),\bar s_v},\Fil_!^{h_0+1}(\Psi_{\xi,\mathds 1_v})_{\mathfrak m}$.
Let us denote
\[
K_0 =\ker (\Fil^{h_0+1}_!(\Psi_{\xi,\mathds 1_v}) \onto
P_\xi(h_0+1,\chi_v)(-\sfrac{h_0}{2})).
\]
Recall that, over $\overline \Qm_l$, all the cohomology groups are concentrated
in degree zero. We~then have an exact sequence
\[
0 \to \Gamma_{\xi,\chi_v,\mathfrak m}(I,h_0+1,+) \to
\Gamma_{\xi,\chi_v,\Psi,\mathfrak m}(I,h_0+1)
\to T \to 0,
\]
where $T \hto H^1_{\tor}(\sh_{I^w(h_0),\bar s_v}, K_0)_{\mathfrak m}$.
The statement about $T[l]$ then follows from the previous section.
\end{proof}

\item
Consider now the case $\delta=h_0/2$ and denote as before by the lattice
\[
\Gamma_{\xi,\chi_v,\Psi,\mathfrak m}(I,h_0+1,-)
\]
of $V_{\xi,\chi_v,\mathfrak m}(I,h_0+1)(\sfrac{h_0}{2})$ induced by the free quotient
of $H^0(\sh_{I^w(h_0),\bar s_v},\Psi_{\xi,\mathds 1_v})_{\mathfrak m}$.
Consider $\Fil^1_{!,\chi_v}(\Psi_{\xi,\mathds 1_v}) \harrow \Fil^1_!
(\Psi_{\xi,\mathds 1_v})
\harrow \Psi_{\xi,\mathds 1_v}$.

\begin{remark}
Before $\Fil^1_{!,\chi_v}(\Psi_{\xi,\mathds 1_v})$ was defined as a quotient of
$\Fil^1_!(\Psi_{\xi,\mathds 1_v})$. To separate the characters 
$\chi_v \in \cusp_{-1}(\mathds 1_v)$
in $\Fil^1_!(\Psi_{\xi,\mathds 1_v})$ we start with a filtration of
$j^{=1,*} \Fil^1_!(\Psi_{\xi,\mathds 1_v})$ with
\[
j^{=1,*} \Fil^1_!(\Psi_{\xi,\mathds 1_v}) \otimes_{\overline \Zm_l} \overline \Qm_l
\simeq \bigoplus_{\chi_v \in \cusp_{-1}(\mathds 1_v)} HT_{\xi,\overline \Qm_l}
(\chi_v,\chi_v),
\]
by considering a numbering of $\cusp_{-1}(\mathds 1_v)$. In particular for a fixed
$\chi_v$, when it appears in first or last position, we may have to modify the lattice
of $HT(\chi_v,\chi_v)$. But as, up to homothety, $HT(\chi_v,\chi_v)$ has a unique
stable lattice, we then obtain the same $\Fil^1_{!,\chi_v}(\Psi_{\xi,\mathds 1_v})$.
\end{remark}

\begin{prop} \label{prop-lattice-noniso}
The lattices $\Gamma_{\xi,\chi_v,\Psi,\mathfrak m} (I,h_0+1,-)$ and
$\Gamma_{\xi,\chi_v,\Psi,\mathfrak m} (I,h_0+1,+)$ are not isomorphic.
\end{prop}

\begin{proof}
Recall that the $\mathfrak m$-cohomology of the perverse sheaves
$P_\xi(t,\chi_v)(\psfrac{t-1}{2})$
for $t \geq h_0+2$, does not have nontrivial invariant vectors under
$\Iw_v(d-h_0,1,\dots,1)$, then
\[
H^0(\sh_{I^w(h_0),\bar s_v},P_\xi(h_0+1,\chi_v)_{\overline \Zm_l}(\sfrac{h_0}{2}))_{\mathfrak m}
\hto H^0(\sh_{I^w(h_0),\bar s_v},\Fil^1_!(\Psi_{\xi,\mathds 1_v}))_{\mathfrak m},
\]
and so the lattice $\Gamma_{\xi,\chi_v,!,\mathfrak m}(I,h_0+1,1)$
is such that we have an exact sequence
\begin{equation} \label{eq-lattice2}
0 \to \Gamma_{\xi,\chi_v,\mathfrak m}(I,h_0+1)
\to \Gamma_{\xi,\chi_v,!,\mathfrak m}
(I,h_0+1,1) \to T \to 0,
\end{equation}
where $T \neq (0)$ is such that $T[l]$, as an $\overline \Fm_l$-representation of
$\GL_d(F_w)$, does not contain any
irreducible generic subquotient made of characters.
By construction we~also have
\[
\Gamma_{\xi,\chi_v,\Psi,\mathfrak m}(I,h_0+1,-)=
\Gamma_{\xi,\chi_v,!,\mathfrak m}(I,h_0+1,1),
\]
so by composing the map of Proposition \ref{prop-lattice-psi}
and \eqref{eq-lattice2}, we obtain an exact sequence
\[
0 \to \Gamma_{\xi,\chi_v,\Psi,\mathfrak m}(I,h_0+1,+) \to
\Gamma_{\xi,\chi_v,\Psi,\mathfrak m}(I,h_0+1,-) \to T_1 \to 0,
\]
where the $l$-torsion $T_1$, as an $\overline \Fm_l$-representation of
$\GL_d(F_w)$, does not contain any
irreducible generic subquotient made of characters. We then conclude from
the fact that the modulo
$l$ reduction of these two lattices admits, as an $\overline \Fm_l$-representation of
$\GL_d(F_w)$, a generic subquotient.
\end{proof}
\end{enumeratea}

\subsection{Global lattices and generic representations}
\label{para-final}

We argue on the set $S_{\mathfrak m}(v)$ of eigenvalues modulo $l$ of
$\overline \rho_{\mathfrak m}(\frob_v)$.
By hypothesis there exists $\lambda:=\chi_v(\varpi_v)$ such that
\[
\{\lambda q_v^{h_0/2},\lambda q_v^{h_0/2-1},\dots,\lambda q_v^{-h_0/2}\}
\subset S_{\mathfrak m}(v),
\]
where $h_0 \geq 1$ is defined above, under the hypothesis that there exists nontrivial torsion. More precisely, $\lambda q_v^{h_0/2}$ is the reduction modulo $l$
of the eigenvalue of $\frob_v$ acting on $P_\xi(\chi_v,h_0+1)_{\overline \Zm_l}(\sfrac{h_0}{2})$
such that the torsion of $H^0(\sh_{I,\bar s_v},P_\xi(\chi_v,h_0)_{\overline \Zm_l}(\sfrac{h_0}{2}))_{\mathfrak m}$
is non zero.

Let us explain the strategy which will be developed in the following.
We first start with $\lambda_0 q_v^{h_0/2} \in S_{\mathfrak m}(v)$ and
we want to prove that $S_{\mathfrak m}(v)$ contains a subset
$\{\lambda_1 q_v^{h_0/2},\lambda_1 q_v^{h_0/2-1},\dots,\lambda_1
q_v^{-h_0/2}\}$ such that
\begin{itemize}
\item there exists $r$ with $0 < r < h_0+1$ with
$\lambda_1 q_v^{h_0/2}=\lambda_0 q_v^{h_0/2+r}$,

\item and $\lambda_1 q_v^{h_0/2}$ is the reduction modulo $l$ of the eigenvalue
of $\frob_v$ acting on $P_\xi(\chi'_v,h_0+1)_{\overline \Zm_l}(\sfrac{h_0}{2})$ such that the torsion of
$H^0(\sh_{I,\bar s_v},P_\xi(\chi'_v,h_0)_{\overline \Zm_l}(\sfrac{h_0}{2}))_{\mathfrak m}$
is non zero.
\end{itemize}
We then obtain another interval inside $S_{\mathfrak m}(v)$
strictly containing the previous one and
we can then play again with $\lambda_1 q_v^{h_0}$ and repeat the above
property. At the end we then obtain the full set $\{\lambda_0 q_v^n\sep n \in \Zm\}$
which is of order the order of $q_v$ modulo $l$. But this order is by hypothesis
strictly greater than $d$ although trivially the set $S_{\mathfrak m}(v)$
of eigenvalues of $\rho(\frob_v)$ is of order $\leq d$.

We now explain how to increase the interval as stated above. To do so,
start first with a classical fact concerning $\overline \Zm_l[G]$-modules, where
$G$ is a group. Let then $\Gamma$, $\Gamma_1$ and $\Gamma_2$ be three
$\overline \Zm_l$-free modules with an action of a group $G$ such that we have an exact sequence
\begin{equation}\label{eq-extension}
0 \to \Gamma_1 \to \Gamma \to \Gamma_2
\to 0,
\end{equation}
which is $G$ equivariant. We then suppose that this extension is split
over $\overline \Qm_l$, \ie\[
\Gamma \otimes_{\overline \Zm_l} \overline \Qm_l
\simeq (\Gamma_1 \otimes_{\overline \Zm_l} \overline \Qm_l) \oplus
(\Gamma_2 \otimes_{\overline \Zm_l} \overline \Qm_l).
\]
Let then set
$\Gamma'_2:=\Gamma \cap (\Gamma_2 \otimes_{\overline \Zm_l} \overline \Qm_l)$
and $\Gamma'_1:=\Gamma/\Gamma_2'$ so that the sequence
\[
0 \to \Gamma'_2 \to \Gamma \to \Gamma'_1
\to 0
\]
is exact.
We then have the following commutative diagram
\begin{equation} \label{eq-prop-extension}
\begin{array}{c}
\xymatrix{
& \Gamma_1 \ar@{^{ (}->}[d] \ar@{=}[r] & \Gamma_1 \ar@{^{ (}->}[d] \\
\Gamma'_2 \ar@{^{ (}->}[r] \ar@{=}[d] & \Gamma \ar@{->>}[r] \ar@{->>}[d] &
\Gamma'_1 \ar@{->>}[d] \\ \Gamma'_2 \ar@{^{ (}->}[r] &
\Gamma_2 \ar@{->>}[r] & T, \\
}
\end{array}
\end{equation}
where $T$ is of torsion and zero if and only if \eqref{eq-extension} is split, \ie $\Gamma'_1=\Gamma_1$ and $\Gamma'_2=\Gamma_2$.

\begin{iremark}
In the following we will consider $\Tm_{I,\xi,\mathfrak m}[\gal_{F,S}]$-free modules.
Note that
when $\Gamma_i \otimes_{\overline \Zm_l} \overline \Fm_l$ for $i=1,2$
are isotypic for the Galois action relatively to a character $\overline \chi_i$
such that $\overline \chi_1 \not \simeq \overline \chi_2$, then from the
previous diagram, we have $T=0$ and
$\Gamma \simeq \Gamma_1 \oplus \Gamma_2$.
\end{iremark}

The idea is to consider two distinct filtrations
$\Gamma:=H^0_{\free}(\sh_{I^w(h_0),\bar s_v},\Psi_{\xi,\mathds 1_v})_{\mathfrak m}$ and
we have to be able to go from one filtration to the other one by using
repeatedly the previous diagram.
The aim of this section is to explain that, upon the hypothesis that the
torsion is non zero, such a process is impossible.

\begin{enumeratea}
\item
First filtration:
recall that
\[
\Gamma \otimes_{\overline \Zm_l} \overline \Qm_l \simeq
\bigoplus_{\widetilde{\mathfrak m}} \Pi_{\widetilde{\mathfrak m}}^I \otimes
\rho_{\widetilde{\mathfrak m}},
\]
so that, by fixing any numbering
\[
\{\widetilde{\mathfrak m} \subset \mathfrak m \sep
\Pi_{\widetilde{\mathfrak m}}^I \neq (0)\}=
\{\widetilde{\mathfrak m}_1,\dots, \widetilde{\mathfrak m}_r\},
\]
we define a filtration $\Fil^{\csbullet}(\Gamma)$ with graded parts
$\gr^k(\Gamma)$ which is a stable lattice of
$\Pi^I_{\widetilde{\mathfrak m}_k} \otimes \rho_{\widetilde{\mathfrak m}_k}$.
With the previous notations, we suppose that
\[
\Pi_{\widetilde{\mathfrak m}_r,v}
\simeq \st_{h_0+1}(\chi_v) \times \chi_{v,1} \times \cdots \times \chi_{v,d-h_0-1}
\]
with $\chi_v \in \cusp_{-1}(\mathds 1_v)$. As the reduction modulo $l$ of
$\rho_{\widetilde{\mathfrak m}_r}$ is irreducible, then
$\gr^r(\Gamma)$ is typical, in the sense of \cite[\S5]{scholze-LT}, \ie\[
\gr^r(\Gamma) \simeq \Gamma_r \otimes \Lambda_r,
\]
where $\Gamma_r$ is a stable lattice of $\Pi_{\widetilde{\mathfrak m}_r}$ on
which the Galois action is trivial and $\Lambda_r$ is a stable lattice of
$\rho_{\widetilde{\mathfrak m}_r}$.

\item
The second filtration of $\Gamma$ is the one induced by the filtration
$\Fil^{\csbullet}_!(\Psi_{\xi,\mathds 1_v})$ and we try, using diagrams like those above,
to end up
to the previous lattice $\gr^r(\Gamma) \simeq \Gamma_r \otimes \Lambda_r$.
As explained in the previous section, $\Gamma_{\xi,\chi_v,\Psi,\mathfrak m}
(I,h_0+1,+)$ is a quotient of~$\Gamma$ and so a quotient of~$\gr^r(\Gamma)$
so that
\[
\Gamma_r \simeq \Gamma_{\xi,\chi_v,\Psi,\mathfrak m} (I,h_0+1,+).
\]
At the opposite, we see, \cf the remark before Proposition \ref{prop-lattice-noniso},
$\Gamma_{\xi,\chi_v,\Psi,\mathfrak m} (I,h_0+1,-)$
as a quotient of
$H^0_{\free}(\sh_{I,\bar s_v},\Fil^1_!(\Psi_{\xi,\mathds 1_v}))_{\mathfrak m}$.
From Proposition \ref{prop-lattice-noniso},
\[
\Gamma_{\xi,\chi_v,\Psi,\mathfrak m} (I,h_0+1,-)\quad\text{and}\quad\Gamma_{\xi,\chi_v,\Psi,\mathfrak m} (I,h_0+1,+)
\]
are not isomorphic.

We can now apply the process of Diagram \eqref{eq-prop-extension} with
the successive irreducible subquotients of $H^0(\sh_{I^w(h_0),\bar s_v},
\Psi_{\xi,\mathds 1_v}/\Fil^1_!(\Psi_{\xi,\mathds 1_v}))_{\mathfrak m}$
until $\Gamma_{\xi,\chi_v,\Psi,\mathfrak m} (I,h_0+1,-)$ is~modified during
the exchange from subspace to quotient. From Proposition \ref{prop-lattice-noniso},
the typical properties coming from the irreducibility of $\rho_{\mathfrak m}$,
we then deduce that there should exist $\chi'_v \in \cusp_{-1}(\varrho)$
and a diagram \eqref{eq-prop-extension} with
\hbox{$\Gamma'_2=\Gamma_{\xi,\chi_v,\Psi,\mathfrak m} (I,h_0+1,-)$}
and $\Gamma'_1$ associated to a subquotient of
\[
H^0\Bigl(\sh_{I,\bar s_v},
P_\xi(t,\chi'_v)_{\overline \Zm_l}\Bigl(\frac{t-1}{2}-\delta\Bigr)\Bigr)_{\mathfrak m}.
\]
Note that,
as this $P_\xi(t,\chi'_v)_{\overline \Zm_l}(\psfrac{t-1}{2}- \delta)$ is a subquotient of some
$\gr^k_!(\Psi_{\xi,\mathds 1_v})$ with $k \geq 2$, then we must have
$2 \leq t \leq h_0+1$ and $\delta>0$.

In particular, from the previous important remark,
we also deduce that
\[
\chi'_v(\psfrac{t-1}{2}-\delta) \equiv
\chi_v(\sfrac{h_0}{2}) \mod l,
\]
so that, by denoting $\lambda_1=\chi'_v(\varpi_v)$,
we can write $\lambda_1 q_v^{h_0/2}=\lambda_0 q_v^{h_0/2+r}$
with $0 < r <h_0+1$.

From Lemma \ref{lem-h0I}, we can repeat the same argument
with $\chi'_v$ in place of $\chi_v$ as it was announced.
\end{enumeratea}

